% Options for packages loaded elsewhere
\PassOptionsToPackage{unicode}{hyperref}
\PassOptionsToPackage{hyphens}{url}
\PassOptionsToPackage{dvipsnames,svgnames,x11names}{xcolor}
\documentclass[
  11pt,
  a4paper,
]{article}
\usepackage{xcolor}
\usepackage[margin=25mm]{geometry}
\usepackage{amsmath,amssymb}
\setcounter{secnumdepth}{-\maxdimen} % remove section numbering
\usepackage{iftex}
\ifPDFTeX
  \usepackage[T1]{fontenc}
  \usepackage[utf8]{inputenc}
  \usepackage{textcomp} % provide euro and other symbols
\else % if luatex or xetex
  \usepackage{unicode-math} % this also loads fontspec
  \defaultfontfeatures{Scale=MatchLowercase}
  \defaultfontfeatures[\rmfamily]{Ligatures=TeX,Scale=1}
\fi
\usepackage{lmodern}
\ifPDFTeX\else
  % xetex/luatex font selection
\fi
% Use upquote if available, for straight quotes in verbatim environments
\IfFileExists{upquote.sty}{\usepackage{upquote}}{}
\IfFileExists{microtype.sty}{% use microtype if available
  \usepackage[]{microtype}
  \UseMicrotypeSet[protrusion]{basicmath} % disable protrusion for tt fonts
}{}
\makeatletter
\@ifundefined{KOMAClassName}{% if non-KOMA class
  \IfFileExists{parskip.sty}{%
    \usepackage{parskip}
  }{% else
    \setlength{\parindent}{0pt}
    \setlength{\parskip}{6pt plus 2pt minus 1pt}}
}{% if KOMA class
  \KOMAoptions{parskip=half}}
\makeatother
\usepackage{longtable,booktabs,array}
\newcounter{none} % for unnumbered tables
\usepackage{calc} % for calculating minipage widths
% Correct order of tables after \paragraph or \subparagraph
\usepackage{etoolbox}
\makeatletter
\patchcmd\longtable{\par}{\if@noskipsec\mbox{}\fi\par}{}{}
\makeatother
% Allow footnotes in longtable head/foot
\IfFileExists{footnotehyper.sty}{\usepackage{footnotehyper}}{\usepackage{footnote}}
\makesavenoteenv{longtable}
\setlength{\emergencystretch}{3em} % prevent overfull lines
\providecommand{\tightlist}{%
  \setlength{\itemsep}{0pt}\setlength{\parskip}{0pt}}
\usepackage{microtype}
\usepackage{needspace}
\usepackage{fvextra}
\DefineVerbatimEnvironment{Highlighting}{Verbatim}{breaklines,commandchars=\\\{\}}
\setlength{\emergencystretch}{3em}
\widowpenalty=10000
\clubpenalty=10000
\allowdisplaybreaks
\usepackage{bookmark}
\IfFileExists{xurl.sty}{\usepackage{xurl}}{} % add URL line breaks if available
\urlstyle{same}
\hypersetup{
  pdftitle={Rank plus Weyl sparsity is at least p+1: an operator form of the Chebotarev-Tao uncertainty principle in prime dimension},
  pdfauthor={Lluis Eriksson (Independent researcher)},
  colorlinks=true,
  linkcolor={blue},
  filecolor={Maroon},
  citecolor={Blue},
  urlcolor={blue},
  pdfcreator={LaTeX via pandoc}}

\title{Rank plus Weyl sparsity is at least p+1: an operator form of the
Chebotarev-Tao uncertainty principle in prime dimension}
\author{Lluis Eriksson (Independent researcher)}
\date{9 September 2026 - Internal revision v004}

\begin{document}
\maketitle

\subsection{Abstract}\label{abstract}

Let \(p\) be prime and let \(W_{(a,\allowbreak{}b)}=X^aZ^b\) be the cyclic Weyl
operators on \(\mathbb C^p\). We prove that every nonzero linear
combination \(C\) of \(\ell\) Weyl operators satisfies
\(\dim\ker C\le\ell-1\), i.e.~\(\operatorname{rank}C+s_W(C)\ge p+1\),
where \(s_W\) counts the nonzero Weyl coefficients. The collinear case
is Tao's uncertainty principle on \(\mathbb Z/p\) (Chebotarev's theorem
on Fourier minors). The general case follows from the
reduction-modulo-\((1-\omega)\) mechanism of Frenkel and Tao after two
further observations: Weyl combinations stay integral under the
non-unimodular Clifford conjugations \(F'Q_s\), which realise all
\(p+1\) line directions of \(\mathbb F_p^2\), and, in the nonvacuous
sparsity range, a pigeonhole over those directions produces a nonzero
fibre sum. The bound is sharp for every \(\ell\le p\) and fails in every
composite dimension. As a self-contained corollary, for one use of the
\(p^2\) Weyl channels the minimal perfect list length over probes of
Schmidt rank \(r\) is exactly \(p-r+1\). An exact certificate refutes a
natural ``line conjecture''.

\subsection{1. Introduction}\label{introduction}

\subsubsection{1.1 The result}\label{the-result}

Fix a prime \(p\), \(\omega=e^{2\pi i/p}\), and the Weyl basis
\(W_{(a,\allowbreak{}b)}=X^aZ^b\) of \(M_p(\mathbb C)\), where \(X\) is the cyclic
shift and \(Z=\operatorname{diag}(\omega^j)\). Every matrix \(C\) has a
unique expansion \(C=\sum_gc_gW_g\) over \(g\in\mathbb F_p^2\); its
\textbf{Weyl sparsity} \(s_W(C)\) is the number of nonzero coefficients.
The main result (Theorem 1, Section 3) is

\[\operatorname{rank}C+s_W(C)\ \ge\ p+1\qquad\text{for every nonzero }C\in M_p(\mathbb C),\]

with equality attained at every sparsity \(1\le\ell\le p\) (Section
5.1). When the support lies on a line of \(\mathbb F_p^2\), \(C\) is (up
to a Weyl translation and Clifford conjugation) a polynomial \(f(Z)\)
with \(\ell\) monomials and the statement reduces to Tao's principle
``an \(\ell\)-term polynomial has at most \(\ell-1\) zeros among the
\(p\)-th roots of unity'', which is Chebotarev's theorem on the minors
of the Fourier matrix. The new content is the non-collinear case. The
proof (Section 4) is short: it is the Frenkel-Tao reduction-modulo-\(p\)
argument together with (i) the observation that the non-unimodular
Clifford operators \(U_s=F'Q_s\) (unnormalised Fourier matrix times a
quadratic phase) conjugate Weyl combinations to Weyl combinations with
the \emph{same} coefficients up to \(p\)-th roots of unity, so
integrality is preserved along all \(p+1\) parallel classes of lines of
\(\mathbb F_p^2\), and (ii) a pigeonhole over the \(p+1\) lines through
a coefficient of minimal valuation, which guarantees a direction in
which the reduced matrix is a nonzero circulant. Theorem 1' is the
sharper valuation-theoretic form. Section 5 shows that the bound is
sharp for every \(\ell\le p\), that every intermediate corank is
attained at exact sparsity \(\ell\), that the multiplicative bound alone
does not imply it (the multiplicative bound
\(\operatorname{rank}C\cdot s_W(C)\ge d\) holds in every dimension), and
exactly where primality enters, matching the composite counterexamples
of {[}0WE{]}. An exact certificate over \(\mathbb Z[\omega_5]\) refutes
the natural refinement ``nullity \(\le\) (largest collinear subset)
\(-1\)'' (Proposition 5).

Section 6 gives a self-contained application to quantum list decoding:
for one use of the \(p^2\) channels \(\operatorname{Ad}(W_g)\) with a
pure probe of Schmidt rank \(r\), the least list length admitting a
perfect decoder, minimised over all such probes, is exactly \(p-r+1\).
The lower bound is Theorem 1 through a sparse-factorisation criterion
(Lemma 6, a Weyl-covariant form of Lemma 9 and Corollary 11 of Johnston,
Lovitz, Russo and Sikora {[}JLRS25{]}); the upper bound is an explicit
decoder for the consecutive probe (Lemma 7, the mechanism of {[}3H0{]}
Theorem 6.1). Both are proved here with their exact hypotheses.

\subsubsection{1.2 Antecedents}\label{antecedents}

\textbf{Vector-level uncertainty principles.} Chebotarev's theorem
(1926) states that every minor of the \(p\times p\) Fourier matrix
\((\omega^{jk})\) is nonzero; see Stevenhagen and Lenstra {[}SL96{]} for
the history and Frenkel {[}Fr04{]} for a half-page proof. Tao {[}Ta05{]}
observed that it is equivalent to
\(|\operatorname{supp}f|+|\operatorname{supp}\hat f|\ge p+1\) for
\(f\ne0\) on \(\mathbb Z/p\), proved it by reduction modulo the prime
\((1-\omega)\) of \(\mathbb Z[\omega]\), and showed the bound sharp.
Frenkel's proof and Tao's are the same mechanism: after reduction,
\(x^p-1=(x-1)^p\), and an \(m\)-term polynomial cannot vanish to order
\(m\) at \(x=1\) because the Euler operator \(x\,\allowbreak{}d/dx\) applied \(j<m\)
times and evaluated at \(1\) yields a Vandermonde system in the
(distinct modulo \(p\)) exponents. The multiplicative bound
\(|\operatorname{supp}f|\,\allowbreak{}|\operatorname{supp}\hat f|\ge|G|\) for every
finite abelian group is due to Donoho and Stark {[}DS89{]}; Meshulam
{[}Me92, Me06{]} refined it for groups with nontrivial subgroups and
extended it to nonabelian groups (rank-support form, {[}WW21, Thm.
3.9{]}); Ghobber and Jaming {[}GJ11{]} generalised Tao's theorem to
pairs of bases of \(\mathbb C^n\). Later variants of Chebotarev's
theorem remain at the level of vectors: functions with symmetry (Garcia,
Karaali and Katz {[}GKK21{]}), cyclic groups of composite order \(pr\)
(Loukaki {[}Lo25{]}), real and finite-field Fourier matrices (Emmrich
and Kunis {[}EK25{]}), and finite-field-valued functions on
\(\mathbb Z/p\) (Quader, Russell and Sundaram {[}QRS19{]}).

\textbf{Time-frequency shifts.} Krahmer, Pfander and Rashkov {[}KPR08{]}
proved uncertainty principles for the short-time Fourier transform on
finite abelian groups; on \(\mathbb Z/p\) their Proposition 4.4 gives
\(\|V_gf\|_0\ge p(p+1)-\|f\|_0\|g\|_0\) when \(\|f\|_0+\|g\|_0>p\), and
\(\|V_gf\|_0\ge p(p+1)-(p+1-\|f\|_0)(p+1-\|g\|_0)\) otherwise and their
Theorem 4.5 gives \(\|f\|_0+\|V_gf\|_0\ge p^2+1\) for a generic window
\(g\). Since \(V_gf(\lambda)\) is, up to a phase, the Weyl coefficient
of the rank-one operator \(f\otimes\bar g\) at \(\lambda\), these are
statements about the Weyl support of \emph{rank-one} operators, stronger
than Theorem 1 in that case and silent for rank \(\ge2\). Lawrence,
Pfander and Walnut {[}LPW05{]} proved that generic Gabor systems in
prime dimension are full spark, and Malikiosis {[}Ma15{]} extended this
to every dimension; these concern linear independence of the
\emph{vectors} \(W_gv\), not the rank of the \emph{operators}
\(\sum_gc_gW_g\).

\textbf{Operator-level antecedents.} A directly related published
inequality between operator rank and Weyl (Pauli) support is
multiplicative: Bu's ``uncertainty principle for Pauli rank'' {[}Bu24,
arXiv v1, Theorem 1, equation (2){]} states, for an \(n\)-qudit
\emph{state} \(\rho\) expanded in the Heisenberg-Weyl basis,
\(\log\operatorname{rank}\rho+\log\chi_P(\rho)\ge n\log d\), with
equality exactly for stabiliser states, where \(\chi_P\) counts the
nonzero Weyl coefficients. In Bu's Section II normalization,
\(\Xi_\rho(g)=\operatorname{Tr}(\rho w(g)^*)\) and
\(\rho=d^{-n}\sum_g\Xi_\rho(g)w(g)\). For \(n=1,\allowbreak{}d=p\), changing \(w(g)\)
to \(X^aZ^b\) changes only nonzero phases and the indexing; our
expansion coefficients are \(p^{-1}\Xi_\rho(g)\) up to those phases.
Thus \(\chi_P(\rho)=s_W(\rho)\), including the identity coefficient, and
exponentiating gives \(\operatorname{rank}\rho\,\allowbreak{}s_W(\rho)\ge p\). This
is the field-free bound \(\operatorname{rank}C\cdot s_W(C)\ge d\) of
Section 5.2 restricted to positive operators (the proof there is the
same Cauchy-Schwarz argument), and it is exactly the kind of bound that
Section 5.2 shows cannot yield the additive law. Meshulam's nonabelian
rank-support inequality, applied to a function on the Heisenberg group
of order \(p^3\) supported on \(\ell\) lifts of the labels, bounds
\(\ell\) times the total rank of its convolution operator, summed over
all irreducible representations; it does not bound the rank of the
single \(p\)-dimensional block \(\sum_gc_gW_g\). Among the sources
consulted, we found no statement of
\(\operatorname{rank}C+s_W(C)\ge p+1\) for general operators, or of the
corresponding Chebotarev-type theorem for the twisted group algebra of
\(\mathbb F_p^2\). This bounded search does not establish priority
(Appendix A).

\textbf{The Weyl records.} Two earlier records of the author study one
use of the \(d^2\) Weyl channels \(\operatorname{Ad}(W_g)\) with an
entangled probe and a final list-valued measurement. Record
ARR-2026-3H0ZKWJMH18MH9FX {[}3H0{]} proves in its Theorem 6.1 that for
the consecutive probe
\(|\Phi_r\rangle=r^{-1/2}\sum_{x<r}|x\rangle|x\rangle\) a perfect
one-use list decoder exists iff \(\ell\ge d-r+1\); its Remark 6.4
exhibits, when \(r\mid d\), an arithmetic-support probe of the same
Schmidt rank with the smaller threshold \(d/r\), and its Corollary 6.5
shows that \(d/r\) is optimal among all probes of Schmidt rank \(r\)
when \(r\mid d\). The optimum over probes in the nondivisible branch
\(r\nmid d\) was left open, pinned between \(\lceil d/r\rceil\) and
\(d-r+1\). Record ARR-2026-0WESHW4YMM9FG8EK {[}0WE{]} reduces the list
problem to a rank problem through its equation (8), the Weyl-covariant
form of the factor-width criterion of {[}JLRS25{]} (re-derived here as
Lemma 6), and proves rank laws for sparsity \(s_W\le3\): Corollary 2.1
(a non-commuting Weyl trinomial in prime dimension has nullity at most
one), Theorem 3 (for \(d=p^k\) the coranks of nonzero matrices with
\(s_W\le3\) are exactly \(p^j\) and \(2p^j\), \(0\le j<k\), for odd
\(p\)), the two-label criterion (rank \(r<d\) admits two labels iff
\((d-r)\mid d\)), and composite-dimension trinomials with large kernels.

Two elementary facts frame what was open. The field-free bound
\(\operatorname{rank}(C)\cdot s_W(C)\ge d\) settles the sparsities
\(\ell\in\{p-1,\allowbreak{}p\}\), and \(\ell\le3\) is settled by {[}0WE{]} Corollary
2.1 together with Chebotarev's theorem for collinear supports. Hence for
\(p\le5\) the law was already a theorem by assembly, and the open cases
were \(4\le\ell\le p-2\), the first being \(p=7\), \(\ell\in\{4,\allowbreak{}5\}\).

\textbf{Calibration.} We describe the contribution as a short,
apparently new corollary of the Frenkel-Tao reduction-mod-\(p\) method;
the new ingredients are the integrality of Weyl combinations under the
non-unimodular Clifford conjugations \(F'Q_s\) and the pigeonhole over
the \(p+1\) directions. The literature searches (above and Appendix A)
found no prior statement of the operator inequality, but a bounded
negative search does not certify priority. Appendix A records how the
manuscript was produced.

\subsection{2. Setting}\label{setting}

Throughout, \(p\) is a prime, \(\omega=e^{2\pi i/p}\), and \(X,\allowbreak{}Z\) act
on \(\mathbb C^p\) by \(X|j\rangle=|j+1\bmod p\rangle\),
\(Z|j\rangle=\omega^j|j\rangle\). They satisfy \(ZX=\omega XZ\). The
cyclic Weyl basis is \(W_{(a,\allowbreak{}b)}=X^aZ^b\), \((a,\allowbreak{}b)\in\mathbb F_p^2\); it
is orthogonal for the trace form,
\(\operatorname{Tr}(W_g^*W_h)=p\,\allowbreak{}\delta_{gh}\), and it is the fixed
basis of {[}0WE{]} and {[}3H0{]} (a tensor product of qubit Pauli bases
is a different setting and is not considered). For
\(C\in M_p(\mathbb C)\) write \(C=\sum_gc_gW_g\); the \textbf{Weyl
support} is \(S(C)=\{g:c_g\ne0\}\) and the \textbf{Weyl sparsity} is
\(s_W(C)=|S(C)|\). Products of Weyl operators are Weyl operators up to
powers of \(\omega\): \(W_gW_h=\omega^{\langle g,\allowbreak{}h\rangle}W_{g+h}\) for
a bilinear form \(\langle\cdot,\allowbreak{}\cdot\rangle\) on \(\mathbb F_p^2\), and
\(W_g^{-1}=W_g^*\) is a phase times \(W_{-g}\). Consequently
\(\operatorname{Tr}(CW_g)=p\,\allowbreak{}\omega^{\langle-g,\allowbreak{}g\rangle}c_{-g}\), so
\(\{g:\operatorname{Tr}(CW_g)\ne0\}=-S(C)\).

\textbf{Symmetries of the problem.} Three kinds of operations preserve
rank and nullity and act on the Weyl support:

\begin{itemize}
\tightlist
\item
  \emph{Translations.} Left multiplication by \(W_h^{-1}\) (invertible)
  maps \(\sum_gc_gW_g\) to \(\sum_gc_g\,\allowbreak{}\phi_{g,\allowbreak{}h}W_{g-h}\) with phases
  \(\phi_{g,\allowbreak{}h}\in\mu_p\): the support is translated by \(-h\), the
  coefficients are multiplied by roots of unity, and the kernel is
  unchanged. Right multiplication translates the support in the same way
  and preserves the rank.
\item
  \emph{Clifford conjugations.} A unitary \(U\) with
  \(UW_gU^{-1}\in\mathbb C\,\allowbreak{}W_{Mg}\) for all \(g\),
  \(M\in SL_2(\mathbb F_p)\), is a Clifford operator; conjugation by it
  permutes the Weyl basis up to phases, maps the support to \(MS\) (a
  \emph{linear} map, fixing the origin), and preserves Weyl sparsity,
  rank and nullity. Every \(M\in SL_2(\mathbb F_p)\) is realised. For
  odd \(p\), we use the following explicit matrices over
  \(\mathbb Q(\omega)\): the unnormalised Fourier matrix
  \(F'=(\omega^{jk})_{j,\allowbreak{}k}\in M_p(\mathbb Z[\omega])\), with
  \(F'^{-1}=p^{-1}F'^*\), and the unitary quadratic phases
  \(Q_s=\operatorname{diag}(\omega^{s\cdot2^{-1}j^2})_j\) for
  \(s\in\mathbb F_p\), where \(2^{-1}\) is the inverse of \(2\) modulo
  \(p\) (\(p\) odd). The Fourier matrix \(F'\), and hence \(F'Q_s\), is
  invertible but not unitary. Their conjugation action is computed in
  Section 4, Step 3.
\item
  \emph{Transposition and Galois action.} Transposition realises a
  determinant \(-1\) label map, and the Galois automorphisms
  \(\omega\mapsto\omega^t\) realise the other determinants (they
  preserve the rank of matrices with algebraic entries; Step 1 of the
  proof shows this suffices).
\end{itemize}

Together these generate the affine group \(\mathrm{AGL}(2,\allowbreak{}p)\) acting on
supports; the maximal nullity over full-support coefficient vectors is
an \(\mathrm{AGL}(2,\allowbreak{}p)\)-invariant of the support (used in Section
7(d)). Note that a line of \(\mathbb F_p^2\) not through the origin can
be moved onto the \(Z\)-axis \(\{(0,\allowbreak{}b)\}\) only by a translation
\emph{followed by} a Clifford conjugation; conjugation alone fixes the
origin.

\textbf{Lines and directions.} The affine plane \(\mathbb F_p^2\) has
\(p+1\) parallel classes of lines (directions), one for each nonzero
linear form \(\lambda:\mathbb F_p^2\to\mathbb F_p\) up to scalars: the
fibres \(\lambda^{-1}(t)\), \(t\in\mathbb F_p\), are the \(p\) lines of
that class. For a subset \(S\) let \(N_\lambda(S)\) be the number of
\(\lambda\)-fibres that meet \(S\). Any two lines through a common point
\(g_0\) with different directions meet only in \(g_0\), so the \(p+1\)
lines through \(g_0\) partition \(\mathbb F_p^2\setminus\{g_0\}\).

\textbf{Valuation rings.} Let \(K\) be a field containing
\(\mathbb Q(\omega)\). A valuation ring \(\mathcal O\subset K\) with
maximal ideal \(\mathfrak m\) and residue field
\(k=\mathcal O/\mathfrak m\) is called \textbf{admissible} if
\(1-\omega\in\mathfrak m\). Since \(p=\prod_{j=1}^{p-1}(1-\omega^j)\)
and each factor is a unit of \(\mathbb Z[\omega]\) times \(1-\omega\),
admissibility forces \(p\in\mathfrak m\), so \(k\) has characteristic
\(p\) and contains \(\mathbb F_p\); moreover
\(\omega\equiv1\pmod{\mathfrak m}\), hence \(\bar Z=I\) and every power
of \(\omega\) reduces to \(1\). Write \(v\) for the valuation and
\(x\mapsto\bar x\) for reduction \(\mathcal O\to k\), extended entrywise
to matrices. Two rank facts are used: for \(A\in M_p(\mathcal O)\),
\(\operatorname{rank}_KA\ge\operatorname{rank}_k\bar A\) (a nonzero
minor of \(\bar A\) lifts to a minor of \(A\) that is a unit of
\(\mathcal O\)), and for \(A\in M_p(K)\),
\(\operatorname{rank}_KA=\operatorname{rank}_{\mathbb C}A\) whenever
\(K\subseteq\mathbb C\) (rank is the size of the largest nonzero minor,
and the minors lie in \(K\)).

\subsection{3. Statement}\label{statement}

\textbf{Theorem 1 (Weyl sparsity-nullity law in prime dimension).} Let
\(p\) be prime, \(S\subseteq\mathbb F_p^2\) with \(|S|=\ell\), and
\(C=\sum_{g\in S}c_gW_g\) with all \(c_g\in\mathbb C\setminus\{0\}\).
Then
\[\dim\ker C\le\ell-1,\qquad\text{equivalently}\qquad\operatorname{rank}C+s_W(C)\ge p+1.\]
The bound is attained for every \(1\le\ell\le p\):
\(C_\ell=\prod_{i=0}^{\ell-2}(Z-\omega^iI)\) has exactly \(\ell\)
nonzero Weyl coefficients and nullity exactly \(\ell-1\). For
\(\ell\ge p+1\) the inequality is vacuous.

\textbf{Theorem 1' (valuation form).} Let \(C\) be as in Theorem 1 with
\(c_g\in K\) for a field
\(\mathbb Q(\omega)\subseteq K\subseteq\mathbb C\), let
\(\mathcal O\subset K\) be an admissible valuation ring, and scale \(C\)
by a coefficient of minimal valuation so that \(\min_gv(c_g)=0\); let
\(\bar c_g\in k\) be the residues. For a nonzero linear form \(\lambda\)
on \(\mathbb F_p^2\) put
\[\gamma^\lambda_t=\sum_{g\in S:\ \lambda(g)=t}\bar c_g\in k\quad(t\in\mathbb F_p),\qquad m_\lambda=\#\{t:\gamma^\lambda_t\ne0\}.\]
Then for every \(\lambda\) with \(m_\lambda\ge1\),
\[\dim\ker C\le m_\lambda-1.\] Some \(\lambda\) with \(m_\lambda\ge1\)
exists whenever \(\ell\le p+1\), and always
\(m_\lambda\le N_\lambda(S)\le\ell\). Consequently
\(\dim\ker C\le\min(m_\lambda-1)\), where the minimum is over all
admissible \(\mathcal O\) (each with its own scaling) and all
\(\lambda\) with \(m_\lambda\ge1\), provided this indexing family is
nonempty. An empty family supplies no additional bound.

Theorem 1' implies Theorem 1 for algebraic coefficients; Step 1 of the
proof shows that this suffices, or one can apply Theorem 1' directly to
\(K=\mathbb Q(\omega)(c_g)\) with a dominating valuation ring supplied
by {[}Stacks, Tag 00IA{]}.

\textbf{Corollary 2 (list decoding; hypotheses and proof in Section 6).}
For one use of the \(p^2\) channels \(\operatorname{Ad}(W_g)\), every
pure probe of Schmidt rank \(r\) admits a perfect \(\ell\)-list decoder
only if \(\ell\ge p-r+1\), and the consecutive probe attains
\(\ell=p-r+1\). Hence \(\ell^{\rm exact}_{\min}(p,\allowbreak{}r)=p-r+1\) for
\(1\le r\le p\).

\subsection{4. Proof of Theorems 1 and
1'}\label{proof-of-theorems-1-and-1}

\textbf{The case \(p=2\).} Theorem 1 is immediate: at \(\ell=1\) the
matrix is invertible; at \(\ell=2\) a nonzero \(2\times2\) matrix has
nullity at most \(1\); and at \(\ell\ge3\) the bound is vacuous. To
prove Theorem 1' as well, extend the scaled coefficient vector by zeros
and write
\[C=\begin{pmatrix}c_{00}+c_{01}&c_{10}-c_{11}\\c_{10}+c_{11}&c_{00}-c_{01}\end{pmatrix},\qquad \det C=c_{00}^{\,2}-c_{01}^{\,2}-c_{10}^{\,2}+c_{11}^{\,2}.\]
In the residue field of characteristic \(2\),
\[\overline{\det C}=\left(\sum_{a,b\in\mathbb F_2}\bar c_{ab}\right)^2.\]
If \(m_\lambda=1\), the total residue sum is the single nonzero fibre
sum, so \(\det C\ne0\) and \(\dim\ker C=0\). If \(m_\lambda=2\),
nonzeroness of \(C\) gives \(\dim\ker C\le1\). Existence of a nonzero
fibre sum for \(\ell\le3\) follows from Step 5, which applies to every
prime. This proves Theorem 1' for \(p=2\).

Henceforth \(p\) is odd. Theorem 1 is vacuous for \(\ell\ge p+1\), so
its algebraic reduction in Step 1 only concerns \(1\le\ell\le p\). For
Theorem 1', the valuation ring is already given: Steps 2-4 and 6 apply
at any sparsity whenever \(m_\lambda\ge1\).

\textbf{Specialization lemma.} Let \(F\subset\mathbb C\) be a number
field and let \(f_1,\allowbreak{}\ldots,\allowbreak{}f_N\in F[x_1,\allowbreak{}\ldots,\allowbreak{}x_s]\) be homogeneous
polynomials of positive degree. If their common zero set contains a
nonzero complex point, it contains a nonzero point over
\(\overline{\mathbb Q}\).

\emph{Proof.} Put \(I=(f_1,\allowbreak{}\ldots,\allowbreak{}f_N)\). If the only common zero over
\(\overline{\mathbb Q}\) were the origin, Hilbert's Nullstellensatz
would give, for each \(j\), an integer \(n_j\ge1\) and an identity
\(x_j^{n_j}=\sum_i h_{ij}f_i\) in \(\overline{\mathbb Q}[x]\). The same
identities hold in \(\mathbb C[x]\), forcing every complex common zero
to be the origin, a contradiction. The conclusion preserves nonzeroness,
not necessarily every nonzero coordinate. \(\square\)

\textbf{Step 1: reduction to algebraic coefficients; the valuation.} Fix
\(S\) with \(1\le|S|=\ell\le p\) and consider the locus
\(B_S=\{c\in\mathbb C^S:\operatorname{rank}C(c)\le p-\ell\}\). Its
defining \((p-\ell+1)\)-minors are homogeneous polynomials of positive
degree with coefficients in
\(\mathbb Z[\omega]\subset\mathbb Q(\omega)\). The specialization lemma
shows that a nonzero complex point of \(B_S\) gives a nonzero algebraic
point \(c'\). Its support \(S'=\{g:c'_g\ne0\}\subseteq S\) is nonempty
and \(\operatorname{rank}C(c')\le p-\ell\le p-|S'|\), so \(c'\) is a
counterexample to Theorem 1 on \(S'\) with all coefficients nonzero and
algebraic. It therefore suffices to prove Theorem 1 for every support
and all nonzero algebraic coefficients; no induction on \(\ell\) is
needed.

So let \(c_g\in\bar{\mathbb Q}\setminus\{0\}\) and
\(K=\mathbb Q(\omega,\allowbreak{}c_g:g\in S)\), a number field. Let \(\mathfrak P\)
be any prime of the ring of integers \(\mathcal O_K\) above \(p\).
Because \((1-\omega)\) is the unique prime of \(\mathbb Z[\omega]\)
above \(p\) (totally ramified),
\(\mathfrak P\cap\mathbb Z[\omega]=(1-\omega)\), so the localisation
\(\mathcal O=(\mathcal O_K)_{\mathfrak P}\) is an admissible discrete
valuation ring with finite residue field \(k\supseteq\mathbb F_p\).

\emph{Alternative route (extension of valuation rings).} One may skip
the Nullstellensatz: for arbitrary nonzero complex \(c_g\) the field
\(K=\mathbb Q(\omega)(c_g:g\in S)\subset\mathbb C\) is finitely
generated over \(\mathbb Q(\omega)\), and the valuation-domination lemma
{[}Stacks, Tag 00IA{]} applies to the discrete valuation ring
\(\mathbb Z[\omega]_{(1-\omega)}\) of \(\mathbb Q(\omega)\) to give a
valuation ring \(\mathcal O\) with fraction field \(K\) dominating it.
Domination means that the maximal ideal contracts to \((1-\omega)\) in
that local subring, so \(\mathcal O\) is admissible. Its value group
need not be \(\mathbb Z\) and \(k\) need not be finite, but the rest of
the proof uses only that \(\mathcal O\) is an admissible valuation ring;
in particular a coefficient of minimal valuation exists because \(S\) is
finite. This is the form in which Theorem 1' is stated.

\textbf{Step 2: scaling and rank comparison.} Choose \(g_0\in S\) with
\(v(c_{g_0})\) minimal and replace \(C\) by \(c_{g_0}^{-1}C\); the
kernel is unchanged, all \(c_g\) now lie in \(\mathcal O\), and
\(\bar c_{g_0}=1\ne0\). Every Weyl matrix has entries in
\(\{0\}\cup\mu_p\subset\mathbb Z[\omega]\subset\mathcal O\), so
\(C\in M_p(\mathcal O)\). By Section 2,
\(\operatorname{rank}_{\mathbb C}C=\operatorname{rank}_KC\ge\operatorname{rank}_k\bar C\),
and the same holds for any \(\mathcal O\)-combination of Weyl matrices.

\textbf{Step 3: the \(p+1\) directions by explicit Clifford
conjugation.} Direct computation gives
\[F'XF'^{-1}=Z,\qquad F'ZF'^{-1}=X^{-1},\qquad Q_sXQ_s^{-1}=\omega^{s\cdot2^{-1}}XZ^s,\qquad Q_sZQ_s^{-1}=Z.\]
(For the first two:
\(F'X|k\rangle=\sum_j\omega^{j(k+1)}|j\rangle=ZF'|k\rangle\) and
\(F'Z|k\rangle=\sum_j\omega^{(j+1)k}|j\rangle=X^{-1}F'|k\rangle\); for
the third:
\(Q_sXQ_s^{-1}|j\rangle=\omega^{s\cdot2^{-1}((j+1)^2-j^2)}|j+1\rangle=\omega^{s\cdot2^{-1}}\omega^{sj}|j+1\rangle\).)
Using \(ZX=\omega XZ\), \((XZ^s)^a=\omega^{s\,\allowbreak{}a(a-1)/2}X^aZ^{sa}\) and
\(Z^aX^{-c}=\omega^{-ac}X^{-c}Z^a\), one obtains for \(U_s:=F'Q_s\) and
every label \((a,\allowbreak{}b)\) the exact identity in \(M_p(\mathbb Q(\omega))\)
\[U_s\,W_{(a,b)}\,U_s^{-1}=\omega^{e(s,a,b)}\,X^{-(sa+b)}Z^{a},\qquad e(s,a,b)=-s\cdot2^{-1}a^2-ab\in\mathbb F_p.\tag{4.1}\]
Indeed
\(Q_sW_{(a,\allowbreak{}b)}Q_s^{-1}=\omega^{s\cdot2^{-1}a}\,\allowbreak{}\omega^{s\,\allowbreak{}a(a-1)/2}X^aZ^{sa+b}=\omega^{s\cdot2^{-1}a^2}X^aZ^{sa+b}\),
and conjugating by \(F'\) gives
\(\omega^{s\cdot2^{-1}a^2}Z^aX^{-(sa+b)}=\omega^{s\cdot2^{-1}a^2-a(sa+b)}X^{-(sa+b)}Z^a\).
All phases are \(p\)-th roots of unity: no \(2p\)-th roots appear
because the basis is \(X^aZ^b\) (not a symmetrised Weyl basis) and
\(2^{-1}\) is taken in \(\mathbb F_p\). The label map
\((a,\allowbreak{}b)\mapsto(-(sa+b),\allowbreak{}a)\) has first coordinate
\(\lambda_s(a,\allowbreak{}b)=-(sa+b)\), whose fibres are the lines of direction
\((1,\allowbreak{}-s)\). Together with \(U_\infty:=I\), whose ``first coordinate'' is
\(\lambda_\infty(a,\allowbreak{}b)=a\) with fibres the lines of direction \((0,\allowbreak{}1)\),
the \(p+1\) operators \(U\in\{I\}\cup\{U_s:s\in\mathbb F_p\}\) realise
every parallel class of lines exactly once.

For each such \(U\) put
\(C_U:=UCU^{-1}=\sum_{g\in S}c_g\,\allowbreak{}\omega^{e_g}\,\allowbreak{}W_{M_Ug}\). Although
\(U_s\) is not unimodular over \(\mathcal O\) (\(\det F'\) has positive
valuation), \(C_U\) is intrinsically an \(\mathcal O\)-combination of
Weyl matrices, hence lies in \(M_p(\mathcal O)\), and
\(\operatorname{rank}_KC_U=\operatorname{rank}_KC\). Identity (4.1) was
verified exactly in \(\mathbb Z[\omega]\) for all \(s\) and all labels,
\(p\in\{3,\allowbreak{}5,\allowbreak{}7,\allowbreak{}11\}\), by three implementations (Section 7); these
computational checks are not independent editorial reports.

\textbf{Step 4: the reduction is a circulant of fibre sums.} Let
\(\lambda\) be the linear form attached to \(U\). Modulo \(\mathfrak m\)
every \(Z^b\) and every phase \(\omega^e\) becomes \(1\), while \(X\) (a
permutation matrix) becomes the cyclic shift \(\bar X\) over \(k\).
Therefore
\[\bar C_U=\sum_{t\in\mathbb F_p}\gamma^\lambda_t\,\bar X^{\,t},\qquad\gamma^\lambda_t=\sum_{g\in S:\ \lambda(g)=t}\bar c_g,\tag{4.2}\]
that is, \(\bar C_U=\bar g_\lambda(\bar X)\) for the polynomial
\(\bar g_\lambda(x)=\sum_t\gamma^\lambda_tx^t\in k[x]\) of degree \(<p\)
(exponents are representatives in \(\{0,\allowbreak{}\dots,\allowbreak{}p-1\}\), which is
legitimate because \(\bar X^p=I\)).

\textbf{Step 5: pigeonhole over the \(p+1\) lines through \(g_0\).} The
\(p+1\) lines through \(g_0\) (one per direction) partition
\(\mathbb F_p^2\setminus\{g_0\}\), so the sets
\((L\setminus\{g_0\})\cap S\) are \(p+1\) pairwise disjoint subsets of
\(S\setminus\{g_0\}\), a set of \(\ell-1\) elements. If \(\ell-1<p+1\),
i.e.~\(\ell\le p+1\), one of them is empty: for the corresponding
direction \(\lambda\) the fibre through \(g_0\) meets \(S\) only in
\(g_0\), and \(\gamma^\lambda_{\lambda(g_0)}=\bar c_{g_0}=1\ne0\). Hence
\(m_\lambda\ge1\) and \(\bar g_\lambda\ne0\). In every direction
\(m_\lambda\le N_\lambda(S)\le\ell\), since a nonzero fibre sum needs a
fibre meeting \(S\).

\textbf{Step 6: rank of a sparse circulant in characteristic \(p\)
(Frenkel's lemma).} Let \(\bar g\in k[x]\) be nonzero of degree \(<p\)
with exactly \(m\) nonzero monomials \(x^{t_1},\allowbreak{}\dots,\allowbreak{}x^{t_m}\), the
\(t_i\) distinct in \(\{0,\allowbreak{}\dots,\allowbreak{}p-1\}\) and hence distinct in
\(\mathbb F_p\subseteq k\). As a \(k[x]\)-module with \(x\) acting as
\(\bar X\), \(k^p\cong k[x]/(x^p-1)=k[x]/(x-1)^p\) (characteristic
\(p\)), and \(\bar g(\bar X)\) is multiplication by \(\bar g\). Write
\(\bar g=(x-1)^ru\) with \(u(1)\ne0\), so \(u\) is a unit modulo
\((x-1)^p\) and \(r<p\); the kernel of multiplication by \((x-1)^r\) on
\(k[x]/(x-1)^p\) is \((x-1)^{p-r}k[x]/(x-1)^p\), of dimension \(r\).
Thus
\[\operatorname{rank}_k\bar g(\bar X)=p-\operatorname{ord}_{x=1}\bar g.\]
\emph{Claim:} \(\operatorname{ord}_{x=1}\bar g\le m-1\). If \(m\ge p\)
this is trivial, because \(\deg\bar g<p\le m\) and \(\bar g\ne0\)
exclude \((x-1)^m\mid\bar g\); so assume \(m\le p-1\) and suppose
\((x-1)^m\mid\bar g\). Let \(\theta=x\,\allowbreak{}d/dx\) (ordinary derivative; no
Hasse derivatives are needed). Since
\(\theta\big((x-1)^ih\big)=x\big(i(x-1)^{i-1}h+(x-1)^ih'\big)\in(x-1)^{i-1}k[x]\)
in every characteristic, \((x-1)^{m-j}\mid\theta^j\bar g\) for
\(0\le j\le m\), hence \((\theta^j\bar g)(1)=0\) for \(j=0,\allowbreak{}\dots,\allowbreak{}m-1\).
But \(\theta^jx^t=t^jx^t\), so these equations read
\(\sum_{i=1}^m\gamma_{t_i}t_i^{\,\allowbreak{}j}=0\) for \(j=0,\allowbreak{}\dots,\allowbreak{}m-1\): a
Vandermonde system in the \(m\) distinct elements
\(t_i\in\mathbb F_p\subseteq k\), which is invertible over any field
containing \(\mathbb F_p\). So all \(\gamma_{t_i}=0\), contradicting
\(\bar g\ne0\). Therefore
\[\operatorname{rank}_k\bar g(\bar X)\ge p-m+1.\tag{4.3}\]

\textbf{Conclusion.} Take any admissible \(\mathcal O\), any direction
\(\lambda\) with \(m_\lambda\ge1\) (Step 5 provides one when
\(\ell\le p+1\)), and the corresponding \(U\). Then \[\begin{aligned}
\operatorname{rank}_{\mathbb C}C
&=\operatorname{rank}_KC=\operatorname{rank}_KC_U\\
&\ge\operatorname{rank}_k\bar C_U=\operatorname{rank}_k\bar g_\lambda(\bar X)\\
&\ge p-m_\lambda+1\ge p-N_\lambda(S)+1\ge p-\ell+1,
\end{aligned}\] which is Theorem 1' and, for \(\ell\le p+1\), the
inequality of Theorem 1. For \(\ell\ge p+1\) the inequality
\(\dim\ker C\le p\le\ell-1\) is vacuous. Sharpness is proved in Section
5.1. \(\square\)

The only place where the choice of \(\mathcal O\) enters is through the
residues \(\bar c_g\); different primes above \(p\) in \(K\) can give
different residue maps and different optimal directions (Section 7, the
split-prime check), which is why Theorem 1' is stated as a minimum over
\((\mathcal O,\allowbreak{}\lambda)\).

\subsection{5. Sharpness and the role of
primality}\label{sharpness-and-the-role-of-primality}

\subsubsection{\texorpdfstring{5.1 Attainment of every corank at exact
sparsity
\(\ell\)}{5.1 Attainment of every corank at exact sparsity \textbackslash ell}}\label{attainment-of-every-corank-at-exact-sparsity-ell}

For \(S\) on the \(Z\)-axis, \(C=f(Z)\) with \(f(z)=\sum_bc_bz^b\), and
\(\dim\ker C=\#\{j:f(\omega^j)=0\}\); Theorem 1 specialises to Tao's
principle ``an \(\ell\)-term polynomial has at most \(\ell-1\) zeros
among the \(p\)-th roots of unity'' (in any non-vertical direction the
\(\ell\) points lie in distinct fibres, so
\(m_\lambda=\#\{b:\bar c_b\ne0\}\le\ell\)). For the family
\(C_\ell=\prod_{i=0}^{\ell-2}(Z-\omega^iI)\), \(2\le\ell\le p\), the
\(q\)-binomial theorem with \(q=\omega\) and \(n=\ell-1\) gives
\[\prod_{i=0}^{n-1}(z-\omega^i)=\sum_{k=0}^{n}(-1)^{n-k}\omega^{\binom{n-k}{2}}\begin{bmatrix}n\\k\end{bmatrix}_\omega z^k,\]
and the Gaussian binomials
\(\begin{bmatrix}n\\k\end{bmatrix}_\omega=\prod_{i=1}^k\frac{1-\omega^{n-i+1}}{1-\omega^i}\)
are nonzero because \(1\le n<p\) keeps every exponent away from
multiples of \(p\). So \(C_\ell\) has exactly \(\ell\) nonzero Weyl
coefficients and exactly \(\ell-1\) zeros
\(\omega^0,\allowbreak{}\dots,\allowbreak{}\omega^{\ell-2}\) among the \(p\)-th roots of unity:
nullity \(\ell-1\). (Exact check in \(\mathbb Z[\omega_7]\) for all
\(\ell\le7\), Section 7.)

More generally, every corank \(k\in\{0,\allowbreak{}\dots,\allowbreak{}\ell-1\}\) is attained at
exact sparsity \(\ell\le p\): take
\(f(z)=\prod_{i<k}(z-\omega^i)\prod_{j<\ell-1-k}(z-\alpha_j)\). Each
coefficient of \(f\) is a polynomial in \(\alpha=(\alpha_j)\) which is
not identically zero (at \(\alpha_j=\omega^{k+j}\) it is a coefficient
of \(C_\ell\)), so for \(\alpha\) outside a proper Zariski-closed set
all \(\ell\) coefficients are nonzero and no \(\alpha_j\) is a \(p\)-th
root of unity; then \(f(Z)\) has sparsity \(\ell\) and nullity exactly
\(k\). Together with Theorem 1 this extends {[}0WE{]} Theorem 3 from
\(s_W\le3\) to all sparsities in prime dimension:

\textbf{Proposition 3.} For \(d=p\) prime and \(1\le\ell\le p\), the set
of ranks of nonzero matrices with \(s_W\le\ell\) is exactly
\(\{p-\ell+1,\allowbreak{}\dots,\allowbreak{}p\}\); there are no forbidden ranks.

\subsubsection{5.2 The field-free bound and its
limits}\label{the-field-free-bound-and-its-limits}

Let \(P\) be the orthogonal projector onto \(\operatorname{ran}C\).
Trace orthogonality gives \(\operatorname{Tr}(W_g^*C)=p\,\allowbreak{}c_g\), and
\(|\operatorname{Tr}(W_g^*C)|=|\langle PW_g,\allowbreak{}C\rangle_F|\le\|PW_g\|_F\|C\|_F=\sqrt{\operatorname{rank}C}\,\allowbreak{}\|C\|_F\)
with \(\|C\|_F^2=p\sum_h|c_h|^2\). Hence
\(p\,\allowbreak{}|c_g|^2\le\operatorname{rank}(C)\sum_h|c_h|^2\) for every \(g\),
and summing over \(g\in S\),
\[\operatorname{rank}(C)\cdot s_W(C)\ge p,\] a Donoho-Stark-type
inequality, valid verbatim in every dimension \(d\) and for every
nonzero matrix. For density matrices in the \(n\)-qudit Heisenberg-Weyl
basis this is Bu's uncertainty principle for Pauli rank {[}Bu24, arXiv
v1, Theorem 1, equation (2){]} (with the same Cauchy-Schwarz proof, and
equality exactly for stabiliser states); we are not aware of it being
stated for general matrices, but it is elementary. It is sharp at
\(\ell=p\) (the rank-one \(|0\rangle\langle0|=p^{-1}\sum_bZ^b\)) and it
already yields Theorem 1 for \(\ell\in\{p-1,\allowbreak{}p\}\) (rank \(\ge2\), resp.
\(\ge1\)), but for \(4\le\ell\le p-2\) it is far weaker than
\(p-\ell+1\) (for \(p=7\), \(\ell=4\): rank \(\ge2\) versus \(\ge4\)).
Thus this multiplicative inequality alone does not imply the additive
bound. In particular, a proof using only assumptions that hold equally
in all composite dimensions would be insufficient; the present proof
uses primality as described below.

\subsubsection{\texorpdfstring{5.3 Prime powers \(d=p^k\),
\(k\ge2\)}{5.3 Prime powers d=p\^{}k, k\textbackslash ge2}}\label{prime-powers-dpk-kge2}

In \(\mathbb Z[\omega_{p^k}]\) the element \(1-\omega_{p^k}\) still
generates the unique (totally ramified) prime above \(p\), so an
admissible reduction exists and Step 4 still produces a circulant of
size \(p^k\) over a field of characteristic \(p\), with
\(x^{p^k}-1=(x-1)^{p^k}\). What fails is Step 6: the exponents \(t_i\)
are distinct modulo \(p^k\) but not modulo \(p\), and the Vandermonde
argument needs them distinct in the residue field. The example
\(1-x^{p^{k-1}}=(1-x)^{p^{k-1}}\) has \(m=2\) monomials and order
\(p^{k-1}\) at \(x=1\); it is the reduction of \(I-Z^{p^{k-1}}\) (after
a Fourier conjugation), which indeed has \(\ell=2\) and nullity
\(p^{k-1}\). This matches {[}0WE{]} Theorem 3 exactly: for \(d=p^k\) the
coranks attainable with \(s_W\le3\) are \(p^j\) and \(2p^j\),
\(0\le j<k\), and for \(j\ge1\) these exceed \(\ell-1\). The law is thus
a prime-dimension phenomenon and not a prime-power one.

\subsubsection{\texorpdfstring{5.4 Composite \(d\) that is not a prime
power}{5.4 Composite d that is not a prime power}}\label{composite-d-that-is-not-a-prime-power}

If \(d\) has at least two distinct prime factors then \(\Phi_d(1)=1\),
so \(1-\omega_d\) is a unit of \(\mathbb Z[\omega_d]\) and there is no
prime ideal in which \(\omega_d\equiv1\): Step 1 has no starting point.
Correspondingly the law is false. {[}0WE{]} Proposition 4 gives
\(2^{1/4}I+X^3+Z^3\) in \(d=12\) with \(\ell=3\) and nullity \(3\);
{[}0WE{]} eq. (6a) gives \(I+X^3+iX^{-3}Z^{18}\) in \(d=36\) with
\(\ell=3\) and nullity \(6\); and {[}0WE{]} eq. (6c) gives the binomial
\(X+XZ^6=X(I+Z^6)\) in \(d=12\) with \(\ell=2\) and nullity \(6\). The
two failure mechanisms (no admissible prime; non-distinct exponents
modulo \(p\)) are exactly the two places where primality is used in
Section 4.

\subsubsection{5.5 Finer bounds for small supports, and the exact
refutation of the line
conjecture}\label{finer-bounds-for-small-supports-and-the-exact-refutation-of-the-line-conjecture}

Theorem 1 is the universal law, not the finest statement for particular
supports. Theorem 1' already recovers the nullity bound \(1\) for every
non-collinear trinomial, consistent with {[}0WE{]} Corollary 2.1.
Indeed, choose a coefficient of minimal valuation at \(g_0\) and let
\(\lambda\) have fibres parallel to the line through the other two
support points. Since the three points are non-collinear, the fibre
through \(g_0\) is a singleton with nonzero residue, while the remaining
two points lie in one other fibre. Thus \(1\le m_\lambda\le2\), and
Theorem 1' gives \(\dim\ker C\le1\). A further bound, used in Section 7
to certify two support orbits, is:

\textbf{Proposition 4 (line plus one point).} Let \(S\) consist of
\(\ell-1\ge1\) points on an affine line \(L\subset\mathbb F_p^2\) and
exactly one point \(h\notin L\), with all coefficients nonzero. Then
\(\dim\ker C\le1\).

\emph{Proof.} Three rank-preserving normalisations (Section 2,
``Symmetries'') bring \(C\) to a standard form.

\begin{enumerate}
\def\labelenumi{(\roman{enumi})}
\item
  \emph{Translation.} Pick a point \(g_1\in L\) and replace \(C\) by
  \(W_{g_1}^{-1}C\). This does not change the kernel; it translates the
  support by \(-g_1\) and multiplies each coefficient by a root of
  unity, so all coefficients remain nonzero. The line \(L\) becomes
  \(L_0=L-g_1\), a line \emph{through the origin}, i.e.~a
  one-dimensional subspace \(\mathbb F_pu\), and the off-line point
  becomes \(h-g_1\notin L_0\).
\item
  \emph{Clifford conjugation.} Since \(SL_2(\mathbb F_p)\) acts
  transitively on nonzero vectors, choose \(M\in SL_2(\mathbb F_p)\)
  with \(Mu=(0,\allowbreak{}1)\) and a Clifford unitary \(U\) realising \(M\);
  replace \(C\) by \(UCU^{-1}\). The kernel dimension is unchanged, the
  coefficients are multiplied by phases, \(L_0\) becomes the \(Z\)-axis
  \(\{(0,\allowbreak{}b)\}\), and the off-line point becomes \((a,\allowbreak{}b_1)\) with
  \(a\ne0\) (a point of \(\mathbb F_p^2\) not on \(\mathbb F_pu\) is
  mapped to a point not on the \(Z\)-axis). Now \(C=g(Z)+c\,\allowbreak{}X^aZ^{b_1}\)
  with \(g\) a polynomial having exactly \(\ell-1\) nonzero coefficients
  and \(c\ne0\).
\item
  \emph{Right multiplication.} Replace \(C\) by \(CZ^{-b_1}\): the rank
  is unchanged, \(g(Z)Z^{-b_1}=(z^{-b_1}g)(Z)\) is again a polynomial in
  \(Z\) with exactly \(\ell-1\) nonzero coefficients (multiplication by
  \(z^{-b_1}\) permutes the monomials of \(\mathbb C[z]/(z^p-1)\)
  cyclically), and the second term becomes \(cX^a\).
\end{enumerate}

So it suffices to bound the nullity of \(C=g(Z)+cX^a\) with \(a\ne0\),
\(c\ne0\) and \(g\ne0\). In coordinates, \((g(Z)v)_j=g(\omega^j)v_j\)
and \((X^av)_j=v_{j-a}\), so \(Cv=0\) reads
\[g(\omega^j)\,v_j=-c\,v_{j-a}\qquad(j\in\mathbb Z/p).\] Since \(a\ne0\)
and \(p\) is prime, \(j\mapsto j-a\) is a single \(p\)-cycle on
\(\mathbb Z/p\), and because \(c\ne0\) each equation determines
\(v_{j-a}\) from \(v_j\). If \(g(\omega^{j_0})=0\) for some \(j_0\), the
equation at \(j_0\) gives \(v_{j_0-a}=0\), then the equation at
\(j_0-a\) gives \(v_{j_0-2a}=0\), and following the cycle every
coordinate vanishes: nullity \(0\). If \(g(\omega^j)\ne0\) for all
\(j\), each equation also determines \(v_j\) from \(v_{j-a}\); starting
from \(v_0\) one obtains \(v_a,\allowbreak{}v_{2a},\allowbreak{}\dots,\allowbreak{}v_{(p-1)a}\) successively,
and the remaining equation (at \(j=0\), closing the cycle) is one
consistency condition. Hence the kernel is at most one-dimensional.
\(\square\)

It is natural to guess that collinearity is the only source of large
kernels, i.e.~that
\(\dim\ker C\le(\text{size of the largest collinear subset of }S)-1\)
(the ``line conjecture''). This is false, and the counterexample is
exact:

\textbf{Proposition 5 (exact refutation of the line conjecture).} Let
\(p=5\) and
\(S=\{(t,\allowbreak{}t^2):t\in\mathbb F_5\}=\{(0,\allowbreak{}0),\allowbreak{}(1,\allowbreak{}1),\allowbreak{}(2,\allowbreak{}4),\allowbreak{}(3,\allowbreak{}4),\allowbreak{}(4,\allowbreak{}1)\}\), a
parabola with no three points collinear (a line meets \(y=x^2\) in at
most two points). The full-support combination
\[C=I+W_{(1,1)}+\omega^2W_{(2,4)}+\omega^3W_{(3,4)}+W_{(4,1)},\qquad\omega=e^{2\pi i/5},\]
has rank exactly \(3\), i.e.~nullity \(2\), whereas the line conjecture
would allow nullity at most \(1\). Theorem 1 (nullity \(\le4\)) is
respected. In fact the maximal nullity on this support is exactly \(2\):
choose a point \((t_0,\allowbreak{}t_0^2)\) whose coefficient has minimal valuation
and take \(\lambda(a,\allowbreak{}b)=b-2t_0a\). On the parabola,
\(\lambda(t,\allowbreak{}t^2)=(t-t_0)^2-t_0^2\) takes exactly three values, with the
fibre through \(t_0\) a singleton. Hence \(1\le m_\lambda\le3\), and
Theorem 1' gives nullity at most \(2\) for every full-support
coefficient vector. The displayed \(C\) attains it.

The certificate was found through a symmetry: the map
\(\Phi(C)=W_{(1,\allowbreak{}1)}Q_2\,\allowbreak{}C\,\allowbreak{}Q_2^{-1}\) sends \(W_{(t,\allowbreak{}t^2)}\) to
\(\omega^{t^2+t}W_{(t+1,\allowbreak{}(t+1)^2)}\) (by (4.1)-type bookkeeping:
\(Q_2W_{(a,\allowbreak{}b)}Q_2^{-1}=\omega^{a^2}W_{(a,\allowbreak{}2a+b)}\) and
\(W_{(1,\allowbreak{}1)}W_{(a,\allowbreak{}b')}=\omega^aW_{(a+1,\allowbreak{}b'+1)}\)), hence permutes the five
Weyl operators of \(S\) cyclically with total phase \(\omega^{40}=1\).
Its eigenvectors on \(\operatorname{span}\{W_g:g\in S\}\) are
\(C_k=\sum_t\omega^{f_k(t)}W_{(t,\allowbreak{}t^2)}\) with
\(f_k(t)=\sum_{i<t}(i^2+i-k)\), \(k\in\mathbb F_5\); all coefficients
are roots of unity, so the rank can be computed exactly over
\(\mathbb Q(\omega_5)\) through the regular representation over
\(\mathbb Q\) (rank of a \(20\times20\) integer matrix divided by
\(4\)). All five \(C_k\) have rank \(3\); \(C=C_0\) is the case \(k=0\),
with \(f_0=(0,\allowbreak{}0,\allowbreak{}2,\allowbreak{}3,\allowbreak{}0)\). Numerically, \(CC^*\) has eigenvalues
\(10,\allowbreak{}10,\allowbreak{}5,\allowbreak{}0,\allowbreak{}0\). The five exponent vectors, the five \(20\times20\)
matrices and their ranks are stored in
{\small\texttt{repro/\allowbreak{}data/\allowbreak{}parabola\_\allowbreak{}certificate.\allowbreak{}json}}; the workshop that
evaluated version 1 recomputed the rank \(12\) of the \(20\times20\)
matrix independently (Section 7(g)).

\subsection{6. The list-decoding
corollary}\label{the-list-decoding-corollary}

This section is self-contained: it proves the two ingredients (Lemmas 6
and 7) that, together with Theorem 1, give Corollary 2.

\subsubsection{6.1 Setting and hypotheses}\label{setting-and-hypotheses}

\emph{Channels.} Dimension \(d=p\) prime. The \(p^2\) channels are the
unitary channels
\(\operatorname{Ad}(W_g):\sigma\mapsto W_g\sigma W_g^*\),
\(g\in\mathbb F_p^2\); one of them, with unknown label \(g\), is applied
exactly once. Priors are irrelevant for what follows.

\emph{Probes.} A pure probe is a unit vector
\(|\psi\rangle\in\mathbb C^p\otimes\mathbb C^R\) (channel input tensor a
reference of arbitrary finite dimension \(R\)). Using the vectorisation
\(|M\rangle\!\rangle=\sum_{i,\allowbreak{}j}M_{ij}|i\rangle|j\rangle\), which
satisfies \((U\otimes V)|M\rangle\!\rangle=|UMV^{T}\rangle\!\rangle\)
and \(\langle\!\langle M|N\rangle\!\rangle=\operatorname{Tr}(M^*N)\),
write \(|\psi\rangle=|A\rangle\!\rangle\) with
\(A\in\mathbb C^{p\times R}\), \(\operatorname{Tr}(A^*A)=1\). The
reduced state on the channel input is \(\rho=AA^*\), and the Schmidt
rank is
\(r=\operatorname{SR}(\psi)=\operatorname{rank}A=\operatorname{rank}\rho\).
The output states are
\(|\psi_g\rangle=(W_g\otimes I)|A\rangle\!\rangle=|W_gA\rangle\!\rangle\).
All of them lie in \(\mathbb C^p\otimes V\), where
\(V\subseteq\mathbb C^R\) is the row space of \(A\) (of dimension
\(r\)), because \(W_g\) only permutes and rescales the rows of \(A\);
compressing every measurement effect by the projector onto
\(\mathbb C^p\otimes V\) changes no outcome probability, so we may and
do assume \(R=r\) and \(A\in\mathbb C^{p\times r}\) of full column rank
\(r\).

\emph{Decoders.} A one-use strategy consists of the probe, the channel
and a final POVM \(\{M_L\}\) on \(\mathbb C^p\otimes\mathbb C^r\) whose
outcomes are lists \(L\subseteq\mathbb F_p^2\) with \(|L|\le\ell\); zero
effects and unused lists are allowed (``weak exclusion''). The decoder
is \textbf{perfect} if the true label is always in the reported list:
\(\langle\psi_g|M_L|\psi_g\rangle=0\) whenever \(g\notin L\). With a
single channel use, the general parallel, sequential and
indefinite-causal-order testers of Bavaresco, Murao and Quintino
{[}BMQ21, BMQ22{]} all reduce to ``probe, channel, POVM'', so this is
the most general one-use strategy. Let \(\ell_{\min}(\psi)\) be the
least \(\ell\) admitting a perfect decoder for the probe \(\psi\) and
\[\ell^{\rm exact}_{\min}(p,r)=\min_{\operatorname{SR}(\psi)=r}\ell_{\min}(\psi).\]
Since \(\ell_{\min}(\psi)\) depends on \(\psi\) only through \(\rho\)
(purifications of the same \(\rho\) differ by an isometry on the
reference, which can be absorbed in the POVM),
\(\ell^{\rm exact}_{\min}(p,\allowbreak{}r)\) is a minimum over reduced states of
rank \(r\).

\emph{Relation to {[}JLRS25{]}.} An indexed family of (possibly
subnormalised) pure states \(\{|\psi_g\rangle\}\) is called
\(k\)-learnable in {[}JLRS25, Def. 1{]} if there is a POVM indexed by
\(k\)-subsets \(S\) with \(\langle\psi_g|M_S|\psi_g\rangle=0\) whenever
\(g\notin S\); their Lemma 9 states that this holds iff the Gram matrix
\((\langle\psi_g|\psi_h\rangle)\) is \(k\)-incoherent, i.e.~a sum of
rank-one positive matrices \(vv^*\) with \(v\) having at most \(k\)
nonzero entries, and their Corollary 11 identifies the least such \(k\)
with the factor width of the Gram matrix. For the Weyl family,
\(\langle\psi_g|\psi_h\rangle=\operatorname{Tr}(\rho\,\allowbreak{}W_g^*W_h)\) is a
phase times a Weyl coefficient of \(\rho\) at \(h-g\), and the
\(k\)-incoherent decompositions of the Gram matrix correspond to the
sparse factorisations of \(\rho\) in Lemma 6 below. We give the direct
proof in Weyl coordinates, which is short, rather than translating the
phases.

\subsubsection{6.2 The sparse-factorisation
criterion}\label{the-sparse-factorisation-criterion}

\textbf{Lemma 6 (Weyl-covariant factor-width criterion; {[}0WE{]} eq.
(8)).} Under the hypotheses of Section 6.1, a perfect \(\ell\)-list
decoder exists for a probe with reduced state \(\rho\) if and only if
\[\rho=\sum_jC_jC_j^*\qquad\text{with}\qquad s_W(C_j)\le\ell\ \text{ for all }j,\tag{6.1}\]
for finitely many matrices \(C_j\in M_p(\mathbb C)\).

\emph{Proof of necessity.} Let \(\{M_L\}\) be a perfect \(\ell\)-list
POVM on \(\mathbb C^p\otimes\mathbb C^r\). Spectrally decompose each
effect, \(M_L=\sum_i|B_{L,\allowbreak{}i}\rangle\!\rangle\langle\!\langle B_{L,\allowbreak{}i}|\)
with \(B_{L,\allowbreak{}i}\in\mathbb C^{p\times r}\), and index the pairs \((L,\allowbreak{}i)\)
by \(j\), writing \(L_j\) for the list of the \(j\)-th vector.
Perfectness gives, for \(g\notin L_j\),
\[0=\langle\!\langle B_j|W_gA\rangle\!\rangle=\operatorname{Tr}(B_j^*W_gA)=\operatorname{Tr}(AB_j^*\,W_g).\]
By Section 2, \(\{g:\operatorname{Tr}(CW_g)\ne0\}=-S(C)\); hence
\(C_j:=AB_j^*/\sqrt p\) has Weyl support \(S(C_j)\subseteq-L_j\), so
\(s_W(C_j)\le\ell\). Completeness
\(\sum_j|B_j\rangle\!\rangle\langle\!\langle B_j|=I_p\otimes I_r\),
after a partial trace over the first factor (for which
\(\operatorname{Tr}_1|B\rangle\!\rangle\langle\!\langle B|=(B^*B)^{T}\)),
gives \(\sum_jB_j^*B_j=p\,\allowbreak{}I_r\). Therefore
\(\sum_jC_jC_j^*=p^{-1}A\big(\sum_jB_j^*B_j\big)A^*=AA^*=\rho\).

\emph{Proof of sufficiency.} Given (6.1), positivity
(\(C_jC_j^*\le\rho\)) gives
\(\operatorname{ran}C_j\subseteq\operatorname{ran}\rho=\operatorname{ran}A\).
Let \(A^+=(A^*A)^{-1}A^*\) be the left inverse of \(A\) (\(A^+A=I_r\),
\(AA^+\) the projector onto \(\operatorname{ran}A\)) and set
\(B_j^*:=A^+C_j\in\mathbb C^{r\times p}\). Then \(AB_j^*=AA^+C_j=C_j\)
and \(\sum_jB_j^*B_j=A^+\rho\,\allowbreak{}(A^+)^*=A^+AA^*(A^+)^*=I_r\). Define
effects indexed by \((j,\allowbreak{}g)\), \(g\in\mathbb F_p^2\):
\[M_{j,g}=\frac1p\,|W_gB_j\rangle\!\rangle\langle\!\langle W_gB_j|.\]
The Weyl twirl \(\sum_gW_gYW_g^*=p\operatorname{Tr}(Y)\,\allowbreak{}I_p\), applied
to the first tensor factor, gives
\(\sum_gM_{j,\allowbreak{}g}=I_p\otimes(B_j^*B_j)^{T}\), so
\(\sum_{j,\allowbreak{}g}M_{j,\allowbreak{}g}=I_p\otimes I_r\): a POVM. Its outcome probabilities
on the output state of label \(h\) are
\[\langle\psi_h|M_{j,g}|\psi_h\rangle=\frac1p\,|\operatorname{Tr}(B_j^*W_g^*W_hA)|^2=\frac1p\,|\operatorname{Tr}(AB_j^*\,W_g^*W_h)|^2=\frac1p\,|\operatorname{Tr}(C_j\,W_{h-g})|^2\]
(the phase of \(W_g^*W_h\in\mu_pW_{h-g}\) disappears in the modulus),
which vanishes unless \(h-g\in-S(C_j)\), i.e.~unless \(h\in g-S(C_j)\).
Reporting the list \(L_{j,\allowbreak{}g}=g-S(C_j)\), of size \(s_W(C_j)\le\ell\), on
outcome \((j,\allowbreak{}g)\) therefore gives a perfect decoder. \(\square\)

Only the necessity direction is used for the lower bound in Corollary 2;
sufficiency shows that (6.1) is an achievability criterion, and it gives
the upper bound an equivalent form (Remark after Lemma 7).

\subsubsection{6.3 The consecutive-probe
decoder}\label{the-consecutive-probe-decoder}

\textbf{Lemma 7 (mechanism of {[}3H0{]} Theorem 6.1, prime dimension).}
For \(1\le r\le p\), the consecutive probe
\(|\Phi_r\rangle=r^{-1/2}\sum_{x=0}^{r-1}|x\rangle|x\rangle\),
i.e.~\(A=r^{-1/2}\binom{I_r}{0}\), \(\rho=r^{-1}P_r\) with \(P_r\) the
projector onto \(\operatorname{span}\{|0\rangle,\allowbreak{}\dots,\allowbreak{}|r-1\rangle\}\),
admits a perfect list decoder with lists of size \(p-r+1\).

\emph{Proof.} The output states are
\[|\psi_{(a,b)}\rangle=r^{-1/2}\sum_{x=0}^{r-1}\omega^{bx}\,|x+a\rangle|x\rangle .\]
For each \(a\in\mathbb F_p\) let
\(V_a=\operatorname{span}\{|x+a\rangle|x\rangle:0\le x<r\}\) and let
\(J_a:\mathbb C^r\to V_a\) be the isometry
\(|x\rangle\mapsto|x+a\rangle|x\rangle\). The \(p\) subspaces \(V_a\)
are mutually orthogonal (for fixed \(x\) the first factors
\(|x+a\rangle\), \(a\in\mathbb F_p\), are orthonormal), and
\(|\psi_{(a,\allowbreak{}b)}\rangle=J_a\varphi_b\) with the \textbf{phase frame}
\[\varphi_b=r^{-1/2}\,(1,\omega^{b},\omega^{2b},\dots,\omega^{(r-1)b})^{T}\in\mathbb C^r,\qquad b\in\mathbb F_p .\]
Let \(T\in\mathbb C^{r\times p}\) have columns \(\varphi_b\). Then
\(TT^*=\sum_b\varphi_b\varphi_b^*=\frac1r\sum_{x,\allowbreak{}y}\big(\sum_b\omega^{b(x-y)}\big)|x\rangle\langle y|=\frac pr\,\allowbreak{}I_r\):
the frame is tight with bound \(\alpha=p/r\).

For every \((r-1)\)-subset \(E\subseteq\mathbb F_p\) let
\(T_E\in\mathbb C^{r\times(r-1)}\) be the submatrix of columns
\(\varphi_b\), \(b\in E\), and define \(w_E\in\mathbb C^r\) by its
signed cofactors,
\[(w_E)_x=(-1)^x\,\overline{\det T_E^{(x)}}\qquad(0\le x<r),\] where
\(T_E^{(x)}\) is \(T_E\) with row \(x\) deleted (for \(r=1\),
\(E=\emptyset\) and \(w_\emptyset=1\)). Two properties follow from
determinant expansions:

\begin{enumerate}
\def\labelenumi{(\alph{enumi})}
\item
  \emph{Orthogonality.} Laplace expansion along the first column gives
  \[\langle w_E,\varphi_b\rangle=\sum_x(-1)^x\det T_E^{(x)}\,(\varphi_b)_x=\det[\varphi_b\,|\,T_E],\]
  which vanishes for \(b\in E\) (a repeated column). (For \(b\notin E\)
  it is a nonzero Vandermonde-type determinant, since the \(\omega^b\)
  are distinct; this full-spark property is not needed below.)
\item
  \emph{Tightness (Cauchy-Binet).} \[\begin{aligned}
  \sum_E(w_E)_x\overline{(w_E)_y}
  &=(-1)^{x+y}\sum_E\overline{\det T_E^{(x)}}\det T_E^{(y)}\\
  &=(-1)^{x+y}\,\overline{\det\big(T^{(x)}(T^{(y)})^*\big)},
  \end{aligned}\] where \(T^{(x)}\) is \(T\) with row \(x\) deleted;
  \(T^{(x)}(T^{(y)})^*\) is the \((r-1)\times(r-1)\) submatrix of
  \(TT^*=\alpha I_r\) with row \(x\) and column \(y\) deleted, whose
  determinant is \(\alpha^{r-1}\) if \(x=y\) and \(0\) otherwise (a zero
  row appears). Hence \[\sum_{|E|=r-1}w_Ew_E^*=\alpha^{r-1}I_r .\]
\end{enumerate}

Now define the POVM on \(\mathbb C^p\otimes\mathbb C^r\) with outcomes
\((a,\allowbreak{}E)\) and a remainder:
\[M_{a,E}=\alpha^{-(r-1)}\,J_a\,w_Ew_E^*\,J_a^*,\qquad M_{\rm rest}=I-\sum_aJ_aJ_a^*.\]
By (b), \(\sum_EM_{a,\allowbreak{}E}=J_aJ_a^*\) is the projector onto \(V_a\); the
\(V_a\) are orthogonal, so \(M_{\rm rest}\) is the projector onto their
common orthogonal complement and the effects sum to \(I\). On outcome
\((a,\allowbreak{}E)\) report the list \(L_{a,\allowbreak{}E}=\{(a,\allowbreak{}b):b\notin E\}\), of size
\(p-r+1\); on \(M_{\rm rest}\) report any fixed list. Perfectness:
\(\langle\psi_{(a',\allowbreak{}b)}|M_{a,\allowbreak{}E}|\psi_{(a',\allowbreak{}b)}\rangle=0\) for \(a'\ne a\)
(orthogonal sectors), and for \(a'=a\) it equals
\(\alpha^{-(r-1)}|\langle w_E,\allowbreak{}\varphi_b\rangle|^2=0\) when \(b\in E\) by
(a); and \(M_{\rm rest}\) annihilates every output state since they all
lie in \(\bigoplus_aV_a\). So the true label is always in the reported
list. \(\square\)

\emph{Remarks.} (1) At \(r=p\) the construction is dense coding
{[}BW92{]}: \(T\) is \(p^{-1/2}\) times the Fourier matrix, each \(w_E\)
is proportional to the single \(\varphi_b\) with \(b\notin E\), and the
lists have length \(1\). At \(r=1\) it measures the shift \(a\) and
reports all \(p\) phase labels. (2) In the language of Lemma 6, the
decoder corresponds to the factorisation \(\rho=\sum_EC_EC_E^*\) with
\(C_E=(r\alpha^{r-1})^{-1/2}h_E(Z)\), where \(h_E\) is the polynomial of
degree \(<p\) with \(h_E(\omega^j)=(w_E)_j\) for \(j<r\) and
\(h_E(\omega^j)=0\) for \(j\ge r\); by (a) its coefficient vector
vanishes on \(E\), so \(s_W(C_E)=|\mathbb F_p\setminus E|=p-r+1\), and
by (b) the sum of the \(C_EC_E^*\) is \(r^{-1}P_r\). This form, and the
POVM itself, were verified numerically for \(p\in\{3,\allowbreak{}5,\allowbreak{}7\}\) and all
\(r\) (Section 7(g)). (3) The lemma uses only the tightness of the phase
frame, not full spark; full spark is what makes \(p-r+1\)
\emph{necessary} for this particular probe in {[}3H0{]} Theorem 6.1, but
here necessity comes from Theorem 1 and holds for every probe of Schmidt
rank \(r\).

\subsubsection{6.4 Proof of Corollary 2}\label{proof-of-corollary-2}

\emph{Lower bound.} Suppose a perfect \(\ell\)-list decoder exists for a
probe of Schmidt rank \(r\). By Lemma 6 (necessity),
\(\rho=\sum_jC_jC_j^*\) with \(s_W(C_j)\le\ell\); some \(C_j\ne0\)
because \(\rho\ne0\), and
\(\operatorname{ran}C_j\subseteq\operatorname{ran}\rho\) by positivity.
Theorem 1 gives
\[r=\operatorname{rank}\rho\ge\operatorname{rank}C_j\ge p-s_W(C_j)+1\ge p-\ell+1,\]
i.e.~\(\ell\ge p-r+1\).

\emph{Upper bound.} Lemma 7 gives a probe of Schmidt rank \(r\) with a
perfect decoder of list size \(p-r+1\). Hence
\(\ell^{\rm exact}_{\min}(p,\allowbreak{}r)=p-r+1\). The rank-budget version
\(\min_{\operatorname{SR}(\psi)\le r}\ell_{\min}(\psi)\) has the same
value because \(p-r+1\) is decreasing in \(r\). \(\square\)

\textbf{Consistency checks.} At \(r=1\) both bounds give \(\ell=p\), the
support-dimension value \(\lceil p/r\rceil\). At \(r=p\) both give
\(\ell=1\). The two-label criterion of {[}0WE{]} (rank \(r<d\) admits
two labels iff \((d-r)\mid d\)) reads, for \(d=p\), ``iff \(r=p-1\)'',
which agrees with \(p-r+1=2\). The divisor branch of {[}3H0{]} Corollary
6.5 is, in prime dimension, only \(r\in\{1,\allowbreak{}p\}\), where \(d/r\) and
\(p-r+1\) coincide.

\textbf{What changes in the corpus.} {[}3H0{]} closed the
optimised-probe problem only for \(r\mid d\) and left
\(\lceil d/r\rceil\le\min_\psi\ell_{\min}(\psi)\le d-r+1\) otherwise;
{[}0WE{]} added the point \(\ell=2\). In prime dimension the whole curve
is now closed at the \emph{upper} end of that interval: the consecutive
probe is globally optimal for every Schmidt rank, the arithmetic-support
trick of {[}3H0{]} Remark 6.4 has no prime-dimension analogue, and the
support-dimension bound \(\lceil p/r\rceil\) is off by a factor of about
\(r\) for mid-range \(r\). By Proposition 3 there are no forbidden ranks
and no positive-sum surprises of the \(13d/16\) type found in {[}0WE{]}
for \(d=2^k\) (a positive sum can only raise the rank).

\textbf{Non-claims.} Approximate decoding, mixed probes, several channel
uses, composite \(d\), and the Bayes success curves below threshold are
untouched.

\subsection{7. Exact and numerical
verification}\label{exact-and-numerical-verification}

The computations below were reported in the supplied version 2. During
preparation of v003, the workshop reran the scripts for (a), (b), and
(g), inspecting their outputs; (g) also rechecked all five exact ranks
in (c). Historical timings and the other runs below are retained as
reported results, not fresh checks. Scripts and the new execution log
are described in Section 9. Finite tests support the written arguments
but do not prove the general theorems.

\textbf{(a) Every ingredient of the proof, exactly}
({\small\texttt{e2\_\allowbreak{}exact\_\allowbreak{}padic\_\allowbreak{}checks.\allowbreak{}py}}, 52-59 s). Exact arithmetic in
\(\mathbb Z[\omega]=\mathbb Z[x]/\Phi_p\); rank over
\(\mathbb Q(\omega)\) through the regular representation over
\(\mathbb Q\); reduction modulo \((1-\omega)\) by \(\omega\mapsto1\)
followed by reduction modulo \(p\). Identity (4.1) for all \(s\) and all
labels, \(p\in\{3,\allowbreak{}5,\allowbreak{}7,\allowbreak{}11\}\) (\(27+125+343+1331=1826\) cases). Then
\(92\) random instances (\(p=5,\allowbreak{}7,\allowbreak{}11\); \(2\le\ell\le5\); fixed seed) in
four regimes: random support and coefficients; the collinear equality
family; coefficients with positive \((1-\omega)\)-adic valuation forcing
cancellations; integer coefficients whose vertical fibre sums vanish
modulo \(p\). For every instance and every direction: \(UCU^{-1}\)
integral, reduction equal to the fibre-sum circulant (4.2),
\(\operatorname{rank}_{\mathbb Q(\omega)}C\ge\operatorname{rank}_{\mathbb F_p}\)(reduction),
reduced rank \(=p-\operatorname{ord}_1\bar g\) with
\(\operatorname{ord}_1\bar g\le m-1\), and maximum over directions
\(\ge p-\ell+1\). Equality \(\operatorname{rank}C=p-\ell+1\) occurred in
\(23\) instances (the collinear products). The rerun output is identical
to the original modulo timings.

\textbf{(b) Second implementation, including a split prime}
({\small\texttt{review\_\allowbreak{}check\_\allowbreak{}proof.\allowbreak{}py}}, historical runtime 42-47 s; no import
of the first ring implementation). Part A: identity (4.1) with the
explicit phase \(e(s,\allowbreak{}a,\allowbreak{}b)\), exact for \(p\in\{3,\allowbreak{}5,\allowbreak{}7\}\). Part B:
\(p=7\), \(\ell=5\), coefficients in \(A=\mathbb Z[\omega_7,\allowbreak{}\sqrt2]\),
where \((1-\omega)\) \emph{splits} (\(2\equiv3^2\bmod7\)) into two
primes with different residue maps, \(\mathfrak P_3:\sqrt2\mapsto3\) and
\(\mathfrak P_4:\sqrt2\mapsto4\). Fourteen instances (generic;
coefficients with positive valuation at exactly one of the two primes; a
common factor \((1-\omega)^2\) removed by exact scaling; forced
fibre-sum cancellations; the collinear equality family), each tested
against both primes and all \(8\) directions: \(UCU^*\) exactly
divisible by \(7\) in \(A\); reduction equal to the predicted circulant;
\(\mathbb F_7\)-rank \(=7-\operatorname{ord}_1\bar g\) with
\(\operatorname{ord}_1\bar g\le m-1\); exact rank over
\(\mathbb Q(\omega,\allowbreak{}\sqrt2)\) (regular representation of degree \(12\))
at least every reduced rank; maximum over directions \(\ge3\). In
several instances the best direction differs between \(\mathfrak P_3\)
and \(\mathfrak P_4\), illustrating the minimum over
\((\mathcal O,\allowbreak{}\lambda)\) in Theorem 1'. Part C: exhaustive check of Step
6, that every polynomial over \(\mathbb F_p\) of degree \(<p\) with
exactly \(m\) nonzero monomials has \(\operatorname{ord}_{x=1}\le m-1\):
\(p=7\), all \(m\), all exponent sets and all nonzero coefficient
vectors (\(823{,\allowbreak{}}542\) polynomials); \(p=5\), all \(m\); \(p=11\),
\(m\le4\) (\(3.47\) million). Never violated, always attained; the
sanity example \(1-x^7\) over \(\mathbb F_7\) (\(m=2\), order \(7\))
shows where distinctness modulo \(p\) enters. Part D: the family
\(\prod_{i<\ell-1}(Z-\omega^iI)\) in \(\mathbb Z[\omega_7]\) has exactly
\(\ell\) nonzero coefficients and nullity \(\ell-1\) for all
\(\ell\le7\).

\textbf{(c) Exact refutation of the line conjecture}
({\small\texttt{review\_\allowbreak{}line\_\allowbreak{}conjecture\_\allowbreak{}exact.\allowbreak{}py}}, 9-12 s including the
imported numerics). Proposition 5: exact rank \(3\) over
\(\mathbb Q(\omega_5)\) for all five \(C_k\), no three support points
collinear (checked exactly). The script also reproduces, numerically,
the other general-position orbit of \(p=5\), \(\ell=5\) (representative
\(\{(0,\allowbreak{}0),\allowbreak{}(1,\allowbreak{}0),\allowbreak{}(0,\allowbreak{}1),\allowbreak{}(1,\allowbreak{}1),\allowbreak{}(2,\allowbreak{}3)\}\)) reaching
\(\sigma_4\approx3\cdot10^{-16}\) with all \(|c_g|\ge0.29\); its
coefficient ratios were not recognised as low-degree algebraic numbers,
so that orbit remains numerical. The coefficient vector is stored at
full double precision in
{\small\texttt{repro/\allowbreak{}data/\allowbreak{}gp5\_\allowbreak{}numerical\_\allowbreak{}witness.\allowbreak{}json}}, together with a
recomputation of \(\sigma_4\) from the stored numbers.

\textbf{(d) Exhaustive support-orbit enumeration with multistart
numerical searches} ({\small\texttt{e1\_\allowbreak{}orbits\_\allowbreak{}numeric.\allowbreak{}py\ p\ ell\ restarts}}).
The symmetry group of the rank problem is
\(\mathrm{AGL}(2,\allowbreak{}p)=GL_2(\mathbb F_p)\ltimes\mathbb F_p^2\) (Section 2,
``Symmetries''). Orbit representatives are computed by canonical forms
and orbit sizes by stabiliser counts, with
\(\sum|\mathrm{orbit}|=\binom{p^2}{\ell}\) asserted in every run. The
enumeration is exhaustive; the coefficient search is not. For a
representative \(S\) and target nullity \(k\), alternating minimisation
of \(\sum_{i>p-k}\sigma_i(C_S(c))^2\) on \(\|c\|_2=1\) uses
\(20\)-\(60\) random restarts. Write \(c^{(u)}_{S,\allowbreak{}k}\) for the final
iterates actually examined and define the observed residual
\[\widehat\mu_{S,k}:=\min_u\sigma_{p-k+1}(C_S(c^{(u)}_{S,k})),\qquad
\widehat\mu_{p,\ell}:=\min_{S\in\mathcal R}\widehat\mu_{S,\ell},\] where
\(\mathcal R\) is the set of tested representatives. The table reports
floating-point approximations of these finite minima, not certified
global minima or positive lower bounds over all coefficients. In exact
arithmetic on the sampled iterates, the true global minimum is at most
the observed minimum; rounding is not certified either. Numerically
detected nullity uses threshold \(10^{-9}\) and the coefficient sphere
permits zeros, so detection may occur on a sub-support. In the table,
``detected nullity'' is the largest such detection over the
representatives, ``best residual'' is \(\widehat\mu_{p,\allowbreak{}\ell}\), and the
next column identifies where nullity \(\ell-1\) was numerically
detected. The rank symmetry using Galois automorphisms does not by
itself transport normalized singular values to every support in the
orbit.

{\def\LTcaptype{none} % do not increment counter
\begin{longtable}[]{@{}
  >{\raggedright\arraybackslash}p{(\linewidth - 12\tabcolsep) * \real{0.14285714}}
  >{\raggedright\arraybackslash}p{(\linewidth - 12\tabcolsep) * \real{0.14285714}}
  >{\raggedright\arraybackslash}p{(\linewidth - 12\tabcolsep) * \real{0.14285714}}
  >{\raggedright\arraybackslash}p{(\linewidth - 12\tabcolsep) * \real{0.14285714}}
  >{\raggedright\arraybackslash}p{(\linewidth - 12\tabcolsep) * \real{0.14285714}}
  >{\raggedright\arraybackslash}p{(\linewidth - 12\tabcolsep) * \real{0.14285714}}
  >{\raggedright\arraybackslash}p{(\linewidth - 12\tabcolsep) * \real{0.14285714}}@{}}
\toprule\noalign{}
\begin{minipage}[b]{\linewidth}\raggedright
\(p\)
\end{minipage} & \begin{minipage}[b]{\linewidth}\raggedright
\(\ell\)
\end{minipage} & \begin{minipage}[b]{\linewidth}\raggedright
orbits
\end{minipage} & \begin{minipage}[b]{\linewidth}\raggedright
detected nullity
\end{minipage} & \begin{minipage}[b]{\linewidth}\raggedright
best residual
\end{minipage} & \begin{minipage}[b]{\linewidth}\raggedright
detecting \(\ell-1\)
\end{minipage} & \begin{minipage}[b]{\linewidth}\raggedright
time
\end{minipage} \\
\midrule\noalign{}
\endhead
\bottomrule\noalign{}
\endlastfoot
3 & 2 & 1 & 1 & 0.707 & the pair & \textless1 s \\
5 & 2 & 1 & 1 & 0.437 & the pair & \textless1 s \\
5 & 3 & 2 & 2 & 0.407 & collinear only & \textless1 s \\
5 & 4 & 5 & 3 & 0.500 & collinear only & 3 s \\
5 & 5 & 11 & 4 & 1.000 & collinear only & 33 s \\
7 & 2 & 1 & 1 & 0.315 & the pair & \textless1 s \\
7 & 3 & 3 & 2 & 0.218 & 2 collinear orbits & \textless1 s \\
7 & 4 & 8 & 3 & 0.180 & 2 collinear orbits & 17 s \\
7 & 5 & 32 & 4 & 0.241 & collinear only & 186 s \\
11 & 4 & 14 & 3 & 0.027 & 4 collinear orbits & 38 s \\
\end{longtable}
}

No tested search detected nullity \(\ell\). The reported best residuals
are at least \(0.18\) except at \(p=11\) (30 restarts); these are
observations about the finite searches, not lower bounds on all
coefficient vectors. Largest numerically detected nullities per
representative: for \((7,\allowbreak{}4)\) and \((11,\allowbreak{}4)\), collinear \(\to3\);
``three collinear plus one'' \(\to2\) (realised on the collinear
sub-support, since Proposition 4 gives \(\le1\) when the off-line
coefficient is nonzero); general position \(\to1\). For \((7,\allowbreak{}5)\):
collinear \(\to4\); ``four collinear plus one'' \(\to3\); twenty orbits
\(\to2\) (all general-position \(5\)-sets among them); nine orbits
\(\to1\), five of which (``three collinear plus two'') have best
observed \(\sigma_6\) residuals between \(4\cdot10^{-5}\) and
\(2\cdot10^{-3}\), so nullity \(2\) there is numerically undecided (the
relevant best observed \(\sigma_3\) residual was at least \(0.48\) on
those representatives; this numerical observation is not used to prove
the law).

\textbf{(e) Additional exact certification of the statement for \(p=7\),
\(\ell=4\)} ({\small\texttt{e3\_\allowbreak{}groebner\_\allowbreak{}p7\_\allowbreak{}l4.\allowbreak{}py}}), the first case not
covered by {[}0WE{]} and the field-free bound. For each of the \(8\)
orbit representatives, with \(c_0=1\) (legitimate by Step 1), the ideal
generated by \(\Phi_7(w)\) and the \(4\times4\) minors of \(C\) in
\(\mathbb Q[c_1,\allowbreak{}c_2,\allowbreak{}c_3,\allowbreak{}w]\) was tested for a Gröbner basis \(\{1\}\).
Six orbits were certified with at most \(240\) random minors
(\(0.2\)-\(31\) s each); the two ``three collinear plus one'' orbits
were not (a rerun with all \(1225\) minors was stopped after \(25\)
minutes) and are instead certified by Proposition 4 (nullity \(\le1\)).
Given Section 4 this certification is redundant, but it predates the
proof and does not share its logic.

\textbf{(f) Corank attainment and Proposition 4}
({\small\texttt{e4\_\allowbreak{}corank\_\allowbreak{}attainment.\allowbreak{}py}}, 1 s;
{\small\texttt{e5\_\allowbreak{}line\_\allowbreak{}plus\_\allowbreak{}point\_\allowbreak{}and\_\allowbreak{}gp5.\allowbreak{}py}}, 20 s). Every corank
\(k\le\ell-1\) at exact sparsity \(\ell\) for \(p\in\{5,\allowbreak{}7,\allowbreak{}11,\allowbreak{}13\}\) by
the construction of Section 5.1; Proposition 4 on \(68\) random exact
instances in \(\mathbb Z[\omega]\), \(p\in\{5,\allowbreak{}7,\allowbreak{}11\}\), all of rank
\(\ge p-1\).

\textbf{(g) Workshop checks and version-2 additions}
({\small\texttt{w1\_\allowbreak{}workshop\_\allowbreak{}and\_\allowbreak{}data.\allowbreak{}py}}, 6 s). The external workshop that
evaluated version 1 (AIRR taller, 2026-09-08, item TALLER-0001) wrote
its own script from the PDF alone: it verified identity (4.1) entrywise
modulo \(p\) for all parameters and labels with \(p\in\{3,\allowbreak{}5,\allowbreak{}7,\allowbreak{}11\}\)
(\(1826\) identities) and recomputed the rank of the \(20\times20\)
regular-representation matrix of the parabola operator as \(12\) (rank
\(3\), nullity \(2\) over \(\mathbb Q(\omega_5)\)). Both numbers agree
with (a)-(c). {\small\texttt{w1\_\allowbreak{}workshop\_\allowbreak{}and\_\allowbreak{}data.\allowbreak{}py}} re-derives them
independently of the author's ring code (integer exponent arithmetic;
rank by exact Fraction elimination and by sympy), writes the data files
of Section 9, and verifies the decoder of Lemma 7 numerically for
\(p\in\{3,\allowbreak{}5,\allowbreak{}7\}\) and all \(1\le r\le p\): tightness \(TT^*=(p/r)I_r\),
the Cauchy-Binet identity \(\sum_Ew_Ew_E^*=(p/r)^{r-1}I_r\),
orthogonality \(\langle w_E,\allowbreak{}\varphi_b\rangle=0\) for \(b\in E\) (and
\(\ne0\) for \(b\notin E\)), completeness of the POVM, zero probability
for every label outside the reported list, and the factorisation
\(\rho=\sum_EC_EC_E^*\) with \(s_W(C_E)=p-r+1\) of Remark (2).

\subsection{8. Open questions}\label{open-questions}

\begin{enumerate}
\def\labelenumi{\arabic{enumi}.}
\tightlist
\item
  \textbf{Exact maximal nullity per support.} For a support \(S\) let
  \(\nu(S)\) be the maximal nullity over full-support coefficient
  vectors, an \(\mathrm{AGL}(2,\allowbreak{}p)\)-invariant. Known: \(\ell-1\) for
  collinear \(S\); \(\le1\) for non-collinear trinomials ({[}0WE{]} Cor.
  2.1) and for ``line plus one point'' (Proposition 4); \(2\) for the
  parabola in \(\mathbb F_5^2\) (Proposition 5). A dimension count (rank
  \(\le p-n\) has codimension \(n^2\) in \(\mathbb P^{\ell-1}\))
  suggests \(\nu(S)\le\lfloor\sqrt{\ell-1}\rfloor\) in general position,
  consistent with Section 7(d) (\(\ell=4\): \(1\); \(\ell=5\): \(2\)),
  but this is a heuristic, not a conjecture we can support beyond
  \(p\le11\).
\item
  \textbf{A Newton-polygon refinement.} Theorem 1' is a first-order
  statement (residues modulo \(\mathfrak m\)). The first-order bound
  already settles non-collinear trinomials and the parabola of
  Proposition 5 (Section 5.5). Can higher-order information sharpen the
  bound for other supports and determine more values of \(\nu(S)\)?
\item
  \textbf{Prime powers.} For \(d=p^k\) the reduction step survives and
  only the Vandermonde step fails (Section 5.3). The \(s_W\le3\) data of
  {[}0WE{]} Theorem 3 (coranks \(p^j\), \(2p^j\)) are consistent with
  \(\dim\ker C\le(\ell-1)p^{k-1}\), which for collinear supports is the
  question of how many \(p^k\)-th roots of unity an \(\ell\)-term
  polynomial can annihilate; we have not investigated it.
\item
  \textbf{Real and orthogonal analogues.} Chebotarev-type statements
  exist for real Fourier matrices and over finite fields {[}GKK21,
  EK25{]}; is there a rank-sparsity law for real Clifford-covariant
  operator bases, or for the Weyl operators of other finite groups (for
  instance groups of order \(pq\), where Meshulam's inequality replaces
  Tao's)?
\item
  \textbf{Quantitative versions.} The finite-search residuals of Section
  7(d) motivate explicit quantitative bounds on
  \(\sigma_{p-\ell+1}(C)/\|c\|\) in terms of \(p\) and \(\ell\). Theorem
  1 and compactness give a qualitative positive minimum for fixed
  \(p,\allowbreak{}\ell\) in the nonvacuous range, including sub-supports, but do not
  identify it with an experimental residual.
\end{enumerate}

\subsection{9. Reproducibility}\label{reproducibility}

The companion source package contains the manuscript, the received
verification scripts in {\small\texttt{repro/\allowbreak{}}}, and fresh v003 execution logs
in {\small\texttt{verification/\allowbreak{}}}. The final workshop run used Python 3.12,
numpy 2.5.1, scipy 1.18 and sympy 1.14 on Windows. The scripts
{\small\texttt{e2\_\allowbreak{}exact\_\allowbreak{}padic\_\allowbreak{}checks.\allowbreak{}py}}, {\small\texttt{review\_\allowbreak{}check\_\allowbreak{}proof.\allowbreak{}py}}
and {\small\texttt{w1\_\allowbreak{}workshop\_\allowbreak{}and\_\allowbreak{}data.\allowbreak{}py}} were inspected and rerun. Their
assertions passed; the logs and hashes identify the exact files and
distinguish historical outputs from this run. The workshop's
{\small\texttt{final\_\allowbreak{}revision\_\allowbreak{}checks.\allowbreak{}py}} verifies the two repaired
normalization identities and the new residue arguments for \(p=2\),
trinomials and the parabola. These finite checks are not a substitute
for the proofs in Sections 4-6.

The data files in {\small\texttt{repro/\allowbreak{}data/\allowbreak{}}} include the five exact parabola
certificates, full-precision coefficients of a separate numerical
witness, and the decoder checks. The slower orbit search and
Gröbner-basis experiments were not repeated for v003; Section 7
explicitly identifies them as historical. No source text from the cited
literature is redistributed. The source package is prepared as a
supplement; its inclusion in an eventual public record is subject to the
author's approval and repository requirements.

\subsection{Changes in internal revision v004 (9 September
2026)}\label{changes-in-internal-revision-v004-9-september-2026}

\begin{enumerate}
\def\labelenumi{\arabic{enumi}.}
\tightlist
\item
  Replaced global-minimum/floor language by explicitly defined
  finite-search residuals and a numerical nullity threshold; unchanged
  historical values are not claimed as certified bounds.
\item
  Added the exact valuation-domination reference and stated how
  domination gives the required contraction.
\item
  Updated the author-editor disclosure to the publicly deployed
  AIRR-FOUNDER-1.0 rule; no acceptance is claimed.
\end{enumerate}

Added a standalone specialization lemma, matched the
Pauli-characteristic-function normalization to Bu's exact v1 theorem and
equation, and softened the literature-search wording.

The earlier deposited v003 and its model reports are retained. This
revision responds to their findings; it is not a final acceptance.

\subsection{Appendix A. AI assistance, provenance and
conflicts}\label{appendix-a.-ai-assistance-provenance-and-conflicts}

The author-supplied drafts attribute substantial use of generative AI to
development of the theorem, proof, scripts and exposition under the
author's direction. The supplied history names Anthropic agents; that
history and the exact historical model designation have not been
independently authenticated by this workshop. During preparation of the
present revision, OpenAI Codex assisted with mathematical review, local
corrections, additional written arguments, execution of verification
scripts, bibliographic checking and document preparation. The author
remains responsible for the scientific content and must approve this
exact revision before submission.

This preparation is AI-assisted manuscript work. It is not a human peer
review, an independent editorial decision, or a formal AIRR
model-assessment report. Historical model verdicts are retained in the
private workshop history and are not presented as editorial
endorsements. The deposited v003 was subsequently assessed in actual
OpenAI Codex calls by gpt-6-astra and gpt-5.6-sol on 9 September 2026.
Their original reports and subsequent clarifications are retained. Both
sets of comments informed this v004 revision; Astra had also
participated in earlier manuscript preparation. A separate reviewing
context does not erase that history. Neither the v003 reports nor this
disclosure constitute approval of the revised v004 artifact.

The author is the founder/operator of AIRR.SCIENCE, the intended
repository; this is a declared editorial conflict. Submission and
acceptance are separate steps. Under the repository's AIRR-FOUNDER-1.0
rule of 9 September 2026, the author may make the human acceptance
decision after two distinct identified model reviews, with his
author-editor role and prior model participation disclosed. This is not
independent human peer review; the final decision remains pending.

The antecedent discussion is based on the sources cited in Section 1.2.
Selected primary statements and bibliographic records were checked
during workshop preparation. Neither the historical searches nor the
present bounded checks establish an exhaustive literature search or
certify priority.

\subsection{References}\label{references}

{[}0WE{]} L. Eriksson, \emph{Three-term Weyl operators: central
multiplicities, rank laws, and exact list decoding}, Archive for
Rigorous Research, ARR-2026-0WESHW4YMM9FG8EK (2026-09-05).
\href{https://airr.science/papers/ARR-2026-0WESHW4YMM9FG8EK/}{Source}.

{[}3H0{]} L. Eriksson, \emph{Strictly Scalable Exterior Decoders for
Quantum Lists: Exact Full-Spark Widths and Fixed-Probe Weyl Bayes
Curves}, Archive for Rigorous Research, ARR-2026-3H0ZKWJMH18MH9FX
(2026-08-24).
\href{https://airr.science/papers/ARR-2026-3H0ZKWJMH18MH9FX/}{Source}.

{[}BMQ21{]} J. Bavaresco, M. Murao, M. T. Quintino, \emph{Strict
hierarchy between parallel, sequential, and indefinite-causal-order
strategies for channel discrimination}, Phys. Rev.~Lett. 127 (2021)
200504. \href{https://doi.org/10.1103/PhysRevLett.127.200504}{Source}.

{[}BMQ22{]} J. Bavaresco, M. Murao, M. T. Quintino, \emph{Unitary
channel discrimination beyond group structures: advantages of sequential
and indefinite-causal-order strategies}, J. Math. Phys. 63 (2022)
042203; arXiv:2105.13369.
\href{https://doi.org/10.1063/5.0075919}{Source}.

{[}Bu24{]} K. Bu, \emph{Extremality of stabilizer states},
arXiv:2403.13632 (2024).
\href{https://arxiv.org/abs/2403.13632}{Source}.

{[}BW92{]} C. H. Bennett, S. J. Wiesner, \emph{Communication via one-
and two-particle operators on Einstein-Podolsky-Rosen states}, Phys.
Rev.~Lett. 69 (1992) 2881-2884.
\href{https://doi.org/10.1103/PhysRevLett.69.2881}{Source}.

{[}DS89{]} D. L. Donoho, P. B. Stark, \emph{Uncertainty principles and
signal recovery}, SIAM J. Appl. Math. 49 (1989) 906-931.
\href{https://doi.org/10.1137/0149053}{Source}.

{[}EK25{]} T. Emmrich, S. Kunis, \emph{Real and finite field versions of
Chebotarëv's theorem}, arXiv:2506.02947 (2025).
\href{https://arxiv.org/abs/2506.02947}{Source}.

{[}Fr04{]} P. E. Frenkel, \emph{Simple proof of Chebotarev's theorem on
roots of unity}, arXiv:math/0312398 (2004).
\href{https://arxiv.org/abs/math/0312398}{Source}.

{[}GJ11{]} S. Ghobber, P. Jaming, \emph{On uncertainty principles in the
finite dimensional setting}, Linear Algebra Appl. 435 (2011) 751-768.
\href{https://arxiv.org/abs/0903.2923}{Source}.

{[}GKK21{]} S. R. Garcia, G. Karaali, D. J. Katz, \emph{An improved
uncertainty principle for functions with symmetry}, J. Algebra 586
(2021) 899-934; arXiv:1807.07648.
\href{https://arxiv.org/abs/1807.07648}{Source}.

{[}JLRS25{]} N. Johnston, B. Lovitz, V. Russo, J. Sikora, \emph{The
complexity of perfect quantum state classification}, arXiv:2510.20789
(2025). \href{https://arxiv.org/abs/2510.20789}{Source}.

{[}KPR08{]} F. Krahmer, G. E. Pfander, P. Rashkov, \emph{Uncertainty in
time-frequency representations on finite abelian groups and
applications}, Appl. Comput. Harmon. Anal. 25 (2008) 209-225;
arXiv:math/0611493. \href{https://arxiv.org/abs/math/0611493}{Source}.

{[}Lo25{]} M. Loukaki, \emph{Chebotarev's theorem for groups of order
\(pq\) and an uncertainty principle}, arXiv:2412.08600 (2024, revised
2025). \href{https://arxiv.org/abs/2412.08600}{Source}.

{[}LPW05{]} J. Lawrence, G. E. Pfander, D. Walnut, \emph{Linear
independence of Gabor systems in finite dimensional vector spaces}, J.
Fourier Anal. Appl. 11 (2005) 715-726.
\href{https://doi.org/10.1007/s00041-005-5017-6}{Source}.

{[}Ma15{]} R.-D. Malikiosis, \emph{A note on Gabor frames in finite
dimensions}, Appl. Comput. Harmon. Anal. 38 (2015) 318-330.
\href{https://doi.org/10.1016/j.acha.2014.06.004}{Source}.

{[}Me92{]} R. Meshulam, \emph{An uncertainty inequality for groups of
order pq}, European J. Combin. 13 (1992) 401-407.
\href{https://doi.org/10.1016/S0195-6698(05)80019-8}{Source}.

{[}Me06{]} R. Meshulam, \emph{An uncertainty inequality for finite
abelian groups}, European J. Combin. 27 (2006) 63-67.
\href{https://arxiv.org/abs/math/0312407}{Source}.

{[}QRS19{]} S. Quader, A. Russell, R. Sundaram, \emph{Small-support
uncertainty principles on \(\mathbb Z/p\) over finite fields},
arXiv:1906.05179 (2019).
\href{https://arxiv.org/abs/1906.05179}{Source}.

{[}SL96{]} P. Stevenhagen, H. W. Lenstra Jr., \emph{Chebotarëv and his
density theorem}, Math. Intelligencer 18 (1996), no. 2, 26-37 (contains
Chebotarev's 1926 theorem on Fourier minors).
\href{https://pub.math.leidenuniv.nl/~lenstrahw/PUBLICATIONS/1996d/art.pdf}{Source}.

{[}Ta05{]} T. Tao, \emph{An uncertainty principle for cyclic groups of
prime order}, Math. Res. Lett. 12 (2005) 121-127; arXiv:math/0308286.
\href{https://arxiv.org/abs/math/0308286}{Source}.

{[}WW21{]} A. Wigderson, Y. Wigderson, \emph{The uncertainty principle:
variations on a theme}, Bull. Amer. Math. Soc. 58 (2021) 225-261;
arXiv:2006.11206. \href{https://arxiv.org/abs/2006.11206}{Source}.

{[}Stacks{]} The Stacks Project Authors, \emph{Valuation rings}, Lemma
10.50.2, Tag 00IA, consulted 9 September 2026.
\href{https://stacks.math.columbia.edu/tag/00IA}{Source}.

\end{document}
